\documentclass[sigconf]{cidr-2027}

\usepackage{microtype}
\usepackage{enumitem}
\usepackage{xspace}
\usepackage{booktabs}
\usepackage{tikz}
\usepackage{colortbl}

\setlist{nosep,leftmargin=1.25em}

\newtheorem{theorem}{Theorem}
\newtheorem{proposition}{Proposition}
\newtheorem{lemma}{Lemma}
\newtheorem{corollary}{Corollary}
\newtheorem{definition}{Definition}
\theoremstyle{remark}
\newtheorem{example}{Example}
\newtheorem{remark}{Remark}

\newcommand{\InputOrd}{\preceq_I}
\newcommand{\OutputOrd}{\preceq_O}
\newcommand{\InputGe}{\succeq_I}
\newcommand{\OutputGe}{\succeq_O}
% Output event emitting #2 after applying the commitment sequence #1
% (\varepsilon = no commitment).
\newcommand{\Out}[2]{\mathsf{output}_{#1}(#2)}
\newcommand{\LiftedOrd}{\trianglelefteq}
% Order vocabulary (role-based names; change the words here).
% Input order:  "j \inrel i" means i \InputOrd j (non-strict); inputs \ingrow;
%   a common \indesc of inputs \inrel each of them (articles: "an \indesc").
% Output order: "u \outrel o" means o \OutputOrd u; a common \outnoun of
%   outputs refines each of them.
\newcommand{\inrel}{expands\xspace}     % relation, 3rd person
\newcommand{\inrelbase}{expand\xspace}  % relation, base form
\newcommand{\inrelpart}{expanding\xspace}% relation, participle
\newcommand{\inrelpast}{expanded\xspace}% relation, past participle
\newcommand{\ingrows}{expands\xspace}        % growth, 3rd person
\newcommand{\ingrow}{expand\xspace}          % growth, base form
\newcommand{\inmore}{more expanded\xspace}    % informal comparative
\newcommand{\indesc}{expansion\xspace}      % noun
\newcommand{\indescs}{expansions\xspace}
\newcommand{\Indesc}{Expansion\xspace}
\newcommand{\outrel}{refines\xspace}         % relation, 3rd person
\newcommand{\outrelbase}{refine\xspace}      % relation, base form
\newcommand{\outrelpart}{refining\xspace}    % relation, participle
\newcommand{\outrelpast}{refined\xspace}     % relation, past participle
\newcommand{\outnoun}{refinement\xspace}     % noun
\newcommand{\outnouns}{refinements\xspace}
\setlength{\emergencystretch}{1.5em}
\newcommand{\Spec}{\mathcal{S}}
\newcommand{\fiber}[1]{S(#1)}
\newcommand{\dog}{\texttt{dog}}
\newcommand{\catv}{\texttt{cat}}
\newcommand{\Sealed}{\mathit{Sealed}}
\newcommand{\Retired}{\mathit{Retired}}
\newcommand{\Rejected}{\mathit{Rejected}}
\newcommand{\retired}[1]{\textcolor{gray}{#1}}
\newcommand{\sh}[1]{\cellcolor{gray!22}#1}
% Full-width strike line through the row just ended; place after \\.
\newcommand{\strikerow}{%
  \noalign{\vskip-0.5\dimexpr\ht\strutbox+\dp\strutbox\relax}%
  \noalign{\vskip-0.5\arrayrulewidth}%
  \arrayrulecolor{black!65}\hline\arrayrulecolor{black}%
  \noalign{\vskip 0.5\dimexpr\ht\strutbox+\dp\strutbox\relax}%
  \noalign{\vskip-0.5\arrayrulewidth}}
% Pointer to the long version (proofs and extended material).
\newcommand{\extended}{the full paper~\cite{hellerstein2026completecalmfull}}
\newcommand{\fullproof}{Proof in the full paper~\cite{hellerstein2026completecalmfull}.}
\newcommand{\proofinfull}{\begin{proof}In the full paper~\cite{hellerstein2026completecalmfull}.\end{proof}}

% ---------- Comment macro; NEVER EVER DELETE!!! ----------
\newcommand{\jmh}[1]{\textcolor{red}{\textit{#1 -- JMH}}}

\usepackage{xcolor}
\definecolor{newblue}{RGB}{0,70,200}
\definecolor{cutred}{RGB}{190,0,0}
\newcommand{\removednote}[1]{{\color{cutred}\footnotesize\,[removed: #1]\,}}
\begin{document}

\title{Complete CALM: A Unified Criterion for Coordination}

\author{Joseph M. Hellerstein}
\affiliation{%
  \institution{UC Berkeley \& Amazon Web Services}
  \city{Berkeley, CA}
  \country{USA}
}

\begin{abstract}
When does a distributed system intrinsically need coordination?  The
literature has lacked a formalism in which the question can be posed in
general.  CALM, CAP, and CRDTs
answer it for one class of programs, consistency models, or design patterns.  Their answers cannot
be compared or combined, and together they leave basic questions open, such as how coordination repairs inconsistency.  We put them on common
ground.  We provide a minimalist relational definition of a specification: a
table that includes, for each input, all outputs a system may produce.  Each
column carries its own partial order, saying when one input \inrel another and
when one output \outrel another.  A specification is \emph{consistent} when its
runs are confluent: runs whose inputs have a common \indesc can always catch up
to a common state.  Coordination shows up as an irrevocable \emph{commitment}
to constrain the system by removing rows from the specification table.

We lift CALM out of its original setting.  \emph{Functional CALM} generalizes
the CALM Theorem from Datalog queries to any function between ordered domains:
such a specification is consistent exactly when it is monotone, and then no
output ever needs coordination.  We extend CAP.  \emph{Complete CAP} covers
any consistency guarantee: a system can keep every allowed output available
and consistent across a partition if and only if its specification passes
Lifted Monotonicity, and Gilbert and Lynch's theorem follows directly.  We also pose and answer two new questions.  \emph{Complete CALM} handles ambiguous specifications
that allow several outputs for the same input, where the test becomes
\emph{Lifted Monotonicity}: as input \ingrows, some output permitted at the
larger input \outrel whatever was emitted at the smaller one.
\emph{Proper Coordination} identifies which stage of a composition must coordinate,
and certifies that the stages behind it need no coordination of their own.  Highly Available Transactions, Invariant
Confluence, and later refinements of CALM all map into the same framework.
\end{abstract}

\begin{CCSXML}
<ccs2012>
<concept>
<concept_id>10003752.10003790.10003795</concept_id>
<concept_desc>Theory of computation~Distributed computing models</concept_desc>
<concept_significance>500</concept_significance>
</concept>
<concept>
<concept_id>10002951.10003317.10003347.10003356</concept_id>
<concept_desc>Information systems~Database transaction processing</concept_desc>
<concept_significance>300</concept_significance>
</concept>
</ccs2012>
\end{CCSXML}

\ccsdesc[500]{Theory of computation~Distributed computing models}
\ccsdesc[300]{Information systems~Database transaction processing}
\keywords{coordination, monotonicity, CALM, CAP, transactions, CRDTs}

\maketitle
\noindent{\color{cutred}\footnotesize\textbf{Removals by category}\par
\begin{tabular}{lrr}category & blocks & words\\\hline
unclassified & 15 & 1885\\
formal definition/notation & 25 & 1366\\
related-work detail & 7 & 1070\\
example & 8 & 690\\
formal caution & 7 & 350\\
explanation & 14 & 350\\
signposting & 17 & 241\\
result & 3 & 185\\
\end{tabular}\par\smallskip\textbf{Largest removals}\par
1. unclassified, 970 words: \textit{The induced specification is an ordered specification, and the input ...}\par
2. formal definition/notation, 523 words: \textit{At an input [math], either no row remains or some ...}\par
3. unclassified, 346 words: \textit{We now make the mapping precise, and then explain how ...}\par
4. related-work detail, 307 words: \textit{An output is a set of public facts, such as ...}\par
5. example, 260 words: \textit{The uniqueness constraint of Bailis et al.'s payroll example shows ...}\par
6. related-work detail, 258 words: \textit{An invariant is a predicate that every state must satisfy.}\par
7. example, 238 words: \textit{A composite can be jointly monotone in three ways.}\par
8. formal definition/notation, 165 words: \textit{In a functional specification this means [math]: the event removes ...}\par
9. related-work detail, 159 words: \textit{Under inclusion of outputs, with inputs only by facts at ...}\par
10. related-work detail, 150 words: \textit{Later refinements of that theorem change the input order or ...}\par
11. formal caution, 135 words: \textit{Seals are similar, with input in place of output, but ...}\par
12. unclassified, 121 words: \textit{As the reading above shows, it implies Relational Monotonicity and ...}\par
}\par\medskip
\noindent{\color{newblue}Blue: new or edited for CIDR.}\quad{\color{cutred}Red: removed from the full version.}\par



\section{Introduction}
\label{sec:intro}

Concurrent systems rely on coordination mechanisms ranging from mutexes and
atomic instructions to database locks and consensus protocols. These mechanisms prevent concurrent actions from producing undesirable
outcomes. They also bring costs: overhead, latency, and
availability in distributed settings. Avoiding unnecessary coordination
is therefore a central goal in distributed systems and databases. This paper
asks what coordination mechanisms fundamentally accomplish, and
when that effect is intrinsically required.

The CALM Conjecture
proposed a crisp, bidirectional answer: consistency can be achieved without
coordination exactly when the specification is
monotone~\cite{hellerstein2010declarative}. Ameloot et al.\ proved it in both
directions, but only for Datalog queries evaluated by relational
transducer networks~\cite{ameloot2013relational,hellerstein2020keeping}. Gilbert and Lynch proved Brewer's CAP conjecture, but only for linearizable
registers~\cite{brewer2000towards,gilbert2002cap}. Both results are
invoked far more broadly: CALM as a general rule that monotone programs
need no coordination, CAP as a general constraint on distributed systems. The folklore outgrew the proofs.

Related results reach similar conclusions from different starting points. Conflict-free Replicated Data Types (CRDTs) use inflationary updates over
join-semilattices to guarantee coordination-free
convergence~\cite{shapiro2011crdt}. Highly Available Transactions (HATs)
identify transactional guarantees that remain available under
partition~\cite{bailis2013hat}. Invariant confluence asks whether
independently reached, invariant-preserving database states can merge
safely~\cite{bailis2014coordination}. Because these results speak different
formal languages, they do not translate directly. Does CALM encompass CAP
(\S\ref{sec:complete-cap})? Are CRDT join laws necessary, or merely one
sufficient discipline for coordination avoidance
(\S\ref{sec:crdt-instance})? Why do some transactional guarantees
stay available under partition while others do not
(\S\ref{sec:hat-instance})? Is invariant confluence a different test from
CALM's (\S\ref{sec:iconfluence-instance})?

We answer such questions with a single criterion. Our foundation is a minimalist
specification model: a relational table pairing input values with
corresponding outputs that are within spec. The model is purely declarative: it
describes what can be observed about a system, without committing to a
language, protocol or placement. The input domain $(I,\InputOrd)$ carries a partial order saying when one input
\inrel another; the output domain $(O,\OutputOrd)$ carries one saying when one
output \outrel another, preserving every claim the earlier one made. \textcolor{newblue}{State-machine-style specifications can be compiled into this form via standard techniques
in formal verification, 
eliding operational details that are not externally observable (\S\ref{sec:scope}).}

\textcolor{newblue}{Simplicity makes the model general: it captures
not only application invariants but also general-purpose consistency definitions
such as linearizability, causal consistency
and transactional isolation levels (\S\ref{sec:complete-cap},
\S\ref{sec:hat-instance}, and \extended).} A specification is \emph{consistent} when the executions it admits are
\emph{confluent}: states reached along different paths can be brought to
a common state.

Our first result, \emph{Functional CALM}
(\S\ref{sec:functional-calm-theorem}), covers specifications that pair each
input with one output: arbitrary functions between ordered domains. It proves
that consistency is equivalent to monotonicity: \(i\InputOrd j\) implies
\(f(i)\OutputOrd f(j)\), so an output emitted at \(i\) remains valid at \(j\). When monotonicity fails, emitting some output requires irrevocably restricting
later results to those that \outrelbase it. We call this restriction a
\emph{commitment}. For a function, every commitment \emph{seals} input,
ruling out later input that could contradict an output, as an election
official does by closing the polls before announcing a result. The functional
case covers relational query languages and the usual CRDT state
specification.

Functions cannot capture settings like leader election or transaction
ordering, where many outputs are within spec for a given input. There,
commitments gain a second job: \emph{choosing} among permitted outputs, as
when a system picks one leader or one serial order. We define \emph{Lifted
Monotonicity}: as input \ingrows, some output permitted at the larger input
\outrel whatever outputs were emitted at the smaller one. \emph{Complete
CALM} proves that consistency is equivalent to Lifted Monotonicity
(\S\ref{sec:complete-calm-theorem}), allowing CALM to formally address
the ambiguous specifications common in transactional and distributed systems.

Section~\ref{sec:distributed-execution} models distributed execution with participants
exchanging messages and defines consistency over their histories; this is shown
to be equivalent to consistency as defined above, so Functional and Complete
CALM apply to distributed executions unchanged. A commitment is
\emph{coordination} when another run could contradict it, by accepting an
input it sealed or emitting an output it chose against; it is \emph{distributed
coordination} when that run is at another participant.

Complete CALM leads to \emph{Complete CAP} (\S\ref{sec:complete-cap}). Given a
consistency guarantee expressed as a specification, Lifted Monotonicity says
which side of CAP it falls on. A specification with Lifted Monotonicity can keep every output it permits
available and confluent during a partition; any other cannot. The Gilbert--Lynch theorem for linearizable
registers follows directly.

This paper makes five contributions:
\begin{enumerate}
\item
\emph{Ordered relational specifications.} Tables over ordered input and
        output domains describe permitted behavior; commitments seal input or choose
        among outputs.
\item
\emph{Consistency as confluence.} One definition covers general
        consistency models and application-specific requirements, and it is
        equivalent to distributed consistency.
\item
\emph{Functional CALM.} For functions between ordered domains,
        consistency is equivalent to monotonicity.
\item
\emph{Complete CALM and CAP.} In general, consistency is equivalent
        to Lifted Monotonicity, and Complete CAP gives the matching
        availability boundary under partition for any declared set of outputs
        that must stay available.
\item
\emph{Composition and Proper Coordination.} Specifications
        compose by equijoin, and we show that in a composition, upstream
        stages can free downstream stages from making commitments (\S\ref{sec:composition}).
\end{enumerate}

Section~\ref{sec:functional-calm} introduces the framework on an election and
proves Functional CALM; Section~\ref{sec:complete-calm} generalizes it to
ambiguous specifications and proves Complete CALM;
Section~\ref{sec:distributed-execution} gives the distributed reading, defines
coordination, and proves Complete CAP. Section~\ref{sec:composition}
composes specifications and characterizes proper coordination. Section~\ref{sec:instances} recovers several earlier results as consequences
of ours and compares other approaches; Section~\ref{sec:discussion} puts
this paper into the broader context of the CALM agenda. \removednote{\textbf{signposting}, 15 words, rank 69; \textit{The appendices give the remaining proofs and further applications; Appendix collects notation and defined terms.}} \textcolor{newblue}{Full proofs and further applications appear in \extended.}


\section{Functional CALM}
\label{sec:functional-calm}

This section focuses on functional specifications, in which each input has a
single specified output. \removednote{\textbf{unclassified}, 29 words, rank 43; \textit{We walk through the framework informally on an election, make it precise, and prove Functional CALM, ...}} \textcolor{newblue}{We walk through the framework informally on an
election (Section~\ref{sec:election-story}), state it precisely
(Section~\ref{sec:framework-brief}), and
prove Functional CALM (Section~\ref{sec:functional-calm-theorem}), identifying
along the way when an output can be safely emitted without commitment.}

\subsection{The framework by example}
\label{sec:election-story}

\begin{figure*}[t]
  \centering
  \small
  \newcommand{\rowhdr}{\toprule $\mathrm{in}$ & $\mathrm{out}$ \\ \midrule}
  \setlength{\tabcolsep}{6pt}
  \begin{tabular}[t]{@{}l c@{}}
    \multicolumn{2}{@{}l}{(a) Initial state, current input $x\cdot\cdot$}\\
    \rowhdr
    \sh{$\mathbf{x\cdot\cdot}$} & \sh{$\bot$} \\
    $xy\cdot$ & \sh{$\bot$} \\
    $xyx$ & \sh{$x$} \\
    $xyy$ & \sh{$y$} \\
    $xx\cdot$ & \sh{$x$} \\
    $xxx$ & \sh{$x$} \\
    $xxy$ & \sh{$x$} \\
    $x\cdot x$ & \sh{$x$} \\
    $x\cdot y$ & \sh{$\bot$} \\
    \bottomrule
    \multicolumn{2}{@{}l}{\footnotesize upper fiber $\{\bot,x,y\}$; undetermined}\\
    \multicolumn{2}{@{}l}{\footnotesize $\bot$ needs no commitment}
  \end{tabular}
  \hfill
  \begin{tabular}[t]{@{}l c@{}}
    \multicolumn{2}{@{}l}{(b) After $\mathsf{input}(xy\cdot)$}\\
    \rowhdr
    \retired{$x\cdot\cdot$} & \retired{$\bot$} \\
    \sh{$\mathbf{xy\cdot}$} & \sh{$\bot$} \\
    $xyx$ & \sh{$x$} \\
    $xyy$ & \sh{$y$} \\
    \retired{$xx\cdot$} & \retired{$x$} \\
    \retired{$xxx$} & \retired{$x$} \\
    \retired{$xxy$} & \retired{$x$} \\
    \retired{$x\cdot x$} & \retired{$x$} \\
    \retired{$x\cdot y$} & \retired{$\bot$} \\
    \bottomrule
    \multicolumn{2}{@{}l}{\footnotesize upper fiber $\{\bot,x,y\}$; undetermined}\\
    \multicolumn{2}{@{}l}{\footnotesize $\bot$ needs no commitment}
  \end{tabular}
  \hfill
  \begin{tabular}[t]{@{}l c@{}}
    \multicolumn{2}{@{}l}{(c) After $\mathsf{input}(xx\cdot)$ instead}\\
    \rowhdr
    \retired{$x\cdot\cdot$} & \retired{$\bot$} \\
    \retired{$xy\cdot$} & \retired{$\bot$} \\
    \retired{$xyx$} & \retired{$x$} \\
    \retired{$xyy$} & \retired{$y$} \\
    \sh{$\mathbf{xx\cdot}$} & \sh{$x$} \\
    $xxx$ & \sh{$x$} \\
    $xxy$ & \sh{$x$} \\
    \retired{$x\cdot x$} & \retired{$x$} \\
    \retired{$x\cdot y$} & \retired{$\bot$} \\
    \bottomrule
    \multicolumn{2}{@{}l}{\footnotesize upper fiber $\{x\}$; determined}\\
    \multicolumn{2}{@{}l}{\footnotesize $x$ needs no commitment}
  \end{tabular}
  \hfill
  \begin{tabular}[t]{@{}l c@{}}
    \multicolumn{2}{@{}l}{(d) (b), then polls closed}\\
    \rowhdr
    \retired{$x\cdot\cdot$} & \retired{$\bot$} \\
    \sh{$\mathbf{xy\cdot}$} & \sh{$\bot$} \\
    \retired{$xyx$} & \retired{$x$} \\ \strikerow
    \retired{$xyy$} & \retired{$y$} \\ \strikerow
    \retired{$xx\cdot$} & \retired{$x$} \\
    \retired{$xxx$} & \retired{$x$} \\
    \retired{$xxy$} & \retired{$x$} \\
    \retired{$x\cdot x$} & \retired{$x$} \\
    \retired{$x\cdot y$} & \retired{$\bot$} \\
    \bottomrule
    \multicolumn{2}{@{}l}{\footnotesize upper fiber $\{\bot\}$; determined}\\
    \multicolumn{2}{@{}l}{\footnotesize $\bot$ is the result}
  \end{tabular}
  \caption{The election specification and three states reached from it, showing rows
  whose input \inrel $x\cdot\cdot$.  Bold: current input; shaded: the current
  input and its upper fiber.  Gray: removed from the remainder table; struck
  through if by a commitment, otherwise retired by an input event.  Panels (b)
  and (c) follow input events.  In (d) the registrar emits $\bot$ and also
  closes the polls, a seal $\mathsf{output}(\bot)$ alone does not require.}
  \label{fig:election-frames}
\end{figure*}

We introduce the framework with a small election, illustrated in
Figure~\ref{fig:election-frames}. Each row of the specification table pairs
the ballots received so far with the permitted result; orders on the columns
describe how inputs \ingrow and which later outputs \outrelbase earlier ones.

Three voters $v_1,v_2,v_3$ each cast at most one immutable ballot for $x$ or
$y$. The output is whichever candidate has two ballots, or $\bot$ while there
is no winner. An input is the ballots received so far, written as a
three-character string with a dot for a ballot not yet received. For brevity,
the figure shows only inputs \inrelpart $x\cdot\cdot$.

The two orders are where a designer makes decisions. The \emph{input order}
records accumulated ballots, one per voter: $xy\cdot$ \inrel $x\cdot\cdot$, while
$xx\cdot$ and $xy\cdot$ are incomparable. The \emph{output order} says which
announcements are safe to build on: $\bot$ claims nothing, so either winner
\outrel it, while $x$ and $y$ are incomparable. In general, the functional table is \emph{monotone}
when
\[
  i\InputOrd j\quad\Longrightarrow\quad f(i)\OutputOrd f(j).
\]
Here, for example, the input sequence $x\cdot\cdot\InputOrd xy\cdot\InputOrd xyx$
monotonically produces the output sequence $\bot\OutputOrd\bot\OutputOrd x$.

\paragraph{Input and determination.}
Each panel captures the specification state: a pair of a \emph{current input}
and a \emph{remainder table} of rows
still possible from that input. When $v_2$'s ballot for $y$ arrives, the event
$\mathsf{input}(xy\cdot)$ produces panel~(b), retiring rows whose inputs no
longer \inrelbase the ballots received, and treating such inputs as illegal. The \emph{upper fiber} is the set of
outputs in the remaining rows at the current input and its \indescs; in panel~(b) it is
$\{\bot,x,y\}$. If $v_2$ instead votes for $x$, as in panel~(c), the upper fiber is $\{x\}$. The state
is \emph{determined}: input alone has resolved the outcome, without commitment.

Panel~(d), by contrast, involves the \emph{system}: the implementation that
receives the ballots and announces the result
(Section~\ref{sec:systems}). In this example, the system acts as a registrar, emitting $\bot$
and closing the polls at
$xy\cdot$ in the same output event. Its irrevocable \emph{seal commitment}
removes the still-possible $xyx$ and $xyy$ rows but leaves the current input
unchanged, and so makes $\bot$ the final result. Emitting $\bot$ by itself
would need no commitment, since every output still possible \outrel it;
closing the polls is a further decision by the system, not something the
specification requires.

\begin{figure}[t]
  \centering
  \small
  \setlength{\tabcolsep}{6pt}
  \begin{minipage}{0.86\columnwidth}
  (e) As (b), all outputs incomparable, $\bot$ emitted\par\smallskip
  \begin{minipage}[t]{0.44\linewidth}
    \begin{tabular*}{\linewidth}[t]{@{}l@{\extracolsep{\fill}}c@{}}
      \toprule
      $\mathrm{in}$ & $\mathrm{out}$ \\ \midrule
      \retired{$x\cdot\cdot$} & \retired{$\bot$} \\
      \sh{$\mathbf{xy\cdot}$} & \sh{$\bot$} \\
      \retired{$xyx$} & \retired{$x$} \\ \strikerow
      \retired{$xyy$} & \retired{$y$} \\ \strikerow
      \retired{$xx\cdot$} & \retired{$x$} \\
    \end{tabular*}
  \end{minipage}\hfill
  \begin{minipage}[t]{0.44\linewidth}
    \begin{tabular*}{\linewidth}[t]{@{}l@{\extracolsep{\fill}}c@{}}
      \toprule
      $\mathrm{in}$ (cont.) & $\mathrm{out}$ (cont.) \\ \midrule
      \retired{$xxx$} & \retired{$x$} \\
      \retired{$xxy$} & \retired{$x$} \\
      \retired{$x\cdot x$} & \retired{$x$} \\
      \retired{$x\cdot y$} & \retired{$\bot$} \\
      \bottomrule
    \end{tabular*}
    {\footnotesize upper fiber $\{\bot\}$; determined}
  \end{minipage}\par\smallskip
  \end{minipage}
  \caption{Panel (b) under the equality order, after $\mathsf{output}(\bot)$, wrapped
  to preserve space.  The
  seal the registrar chose in Figure~\ref{fig:election-frames}(d) is required
  here.}
  \label{fig:nonmonotone-frame}
\end{figure}

\paragraph{Output and consistency.}
In a functional table, an output $o$ is \emph{exposable} when every output
that could occur in future \outrel it---the entire upper fiber of the current input. Hence $\bot$ is exposable
in panels~(a) and~(b), and $x$ in panel~(c). The atomic output event
$\mathsf{output}(o)$ seals any legal input whose output does not \outrelbase
$o$, then emits $o$, in one transition. Emitting $o$ is a promise:
from then on, only rows whose outputs \outrelbase $o$ remain. The emission
matters because the environment is powerful and unpredictable; it may use $o$
when choosing later input.

Figure~\ref{fig:nonmonotone-frame} considers a variant of our example that
changes only the output order: all outputs are incomparable, including $\bot$
for no winner (the \emph{equality order}, under which each element is related only
to itself). Then $\bot$ is incompatible with future winners, so emitting it
requires a seal. That change in the output order is Functional CALM in
miniature: the original order is monotone with respect to input order, so all outputs remain
valid as input grows; under the equality order, some outputs can be emitted
only with a commitment.

The obstruction is a property solely of the table and its orders. Rows $(xy\cdot,\bot)$ and
$(xyx,x)$ have the common input \indesc $xyx$, but under the equality
order no output at $xyx$ \outrel both outputs $\bot$ and $x$. A run that emits
$\bot$ at $xy\cdot$ and one that receives $xyx$ cannot be reconciled at $xyx$.

Input and output events form a transition system, and the specification is
consistent when its runs are confluent. \removednote{\textbf{unclassified}, 15 words, rank 70; \textit{For functional tables this corresponds exactly to monotonicity (Section ).}} \textcolor{newblue}{Section~\ref{sec:functional-calm-theorem} shows that for functional tables
this corresponds exactly to monotonicity.}

\subsection{The framework in brief}
\label{sec:framework-brief}
\label{sec:functional-specifications}
\label{sec:functional-events}
\label{sec:executions}

\textcolor{newblue}{We now state these notions precisely.}

\begin{definition}[Functional ordered specification]
\label{def:functional-specification}
\textcolor{newblue}{A \emph{functional ordered specification} is a tuple
$\Spec=(I,\InputOrd,O,\OutputOrd,S)$,
where $(I,\InputOrd)$ and $(O,\OutputOrd)$ are partially ordered input and
output types, and $S\subseteq I\times O$ is a typed relation with schema
$(\mathrm{in},\mathrm{out})$ in which every
input appears in at least one row ($S$ is \emph{total} on $I$) and at most one
(a functional dependency).} \textcolor{newblue}{The unique output paired with $i$ induces a
function $f:I\to O$.} \removednote{\textbf{formal definition/notation}, 8 words, rank 83; \textit{with [math] We call [math] the input order .}}
\end{definition}

\textcolor{newblue}{Because an output may depend on everything a system has received, we model
inputs as \emph{cumulative}: $i\InputOrd j$ means that $j$ contains at least
the input information in $i$.} \textcolor{newblue}{We say that $j$ \emph{\inrel} $i$, or is an
\emph{\indesc} of $i$.} We call $\OutputOrd$ the \emph{output order}:
$o\OutputOrd u$ means that $u$ preserves every externally visible claim in $o$
and may add more, and we say that $u$ \emph{\outrel} $o$. \removednote{\textbf{formal definition/notation}, 24 words, rank 53; \textit{We use these two verbs only in these senses: inputs , and outputs are .}} \textcolor{newblue}{A \emph{common
\indesc} of inputs is an input that \inrel each of them, and a \emph{common
\outnoun} of outputs is an output that \outrel each of them.} \removednote{\textbf{formal caution}, 63 words, rank 21; \textit{Neither need be unique or least: the input order need not be a lattice, only a ...}} \removednote{\textbf{formal definition/notation}, 12 words, rank 74; \textit{A functional ordered specification [math], with induced function [math], is monotone if}} \removednote{\textbf{explanation}, 11 words, rank 81; \textit{Thus an output on partial input stays valid as input .}} \textcolor{newblue}{Neither need be
unique or least: for example, two independent transactions $a$ and $b$, with schedules
ordered by subsequence, have two incomparable minimal common \indescs, $ab$
and $ba$.}

\begin{definition}[Functional monotonicity]
\label{def:functional-monotonicity}
\textcolor{newblue}{$\Spec$ is \emph{monotone} if}
\begin{equation}
  i\InputOrd j
  \quad\Longrightarrow\quad
  f(i)\OutputOrd f(j).
  \label{eq:functional-monotonicity}
\end{equation}
\end{definition}

\textcolor{newblue}{We introduce some compact notation.} \textcolor{newblue}{For a relation $R\subseteq I\times O$, a predicate in square brackets selects
rows, and commas separate conjuncts: $R[\mathrm{in}\InputGe j,\
\mathrm{out}\OutputGe o]$ holds the rows whose input \inrel $j$ and whose
output \outrel $o$.} \textcolor{newblue}{We say that a row $(k,u)$ is \emph{at} $k$, and $\mathrm{out}(R)$ is
the set of outputs in $R$.} \removednote{\textbf{formal definition/notation}, 88 words, rank 18; \textit{A specification describes every permitted input/output pair.}; labels def:residual-state}

\paragraph{States.}
\textcolor{newblue}{As illustrated by the non-gray text in the panels of Figure~\ref{fig:election-frames}, a \emph{specification state} is a pair $\sigma=(i,R)$ of a \emph{current
input} $i$ and a \emph{remainder table} $R\subseteq S$ that contains a row at
$i$ and only rows at \indescs of $i$.} An input $k$ is \emph{legal} at
$\sigma$ when $R$ contains a row at $k$. Rows are never added during
execution. \removednote{\textbf{formal definition/notation}, 34 words, rank 37; \textit{The largest state at input [math], and the one a system is in after receiving [math] ...}} \removednote{\textbf{explanation}, 12 words, rank 75; \textit{Input events alone move between canonical states, and runs start at one.}} \textcolor{newblue}{The \emph{canonical state at $i$}, the state of a system that has
received $i$ and emitted nothing, retains all rows at $i$ and its \indescs:}
\begin{equation}
  \sigma_S(i)=\bigl(i,\ S[\mathrm{in}\InputGe i]\bigr).
  \label{eq:canonical-state}
\end{equation}
\textcolor{newblue}{The \emph{fiber} of $R$ at $k$ is $R(k)=\mathrm{out}(R[\mathrm{in}=k])$.} \removednote{\textbf{formal definition/notation}, 39 words, rank 32; \textit{A functional fiber has at most one value.}} \removednote{\textbf{unclassified}, 26 words, rank 48; \textit{Two events change state: input events incorporate new information, while output events record promises to the ...}} \textcolor{newblue}{The
\emph{upper fiber} of $(i,R)$ is $\mathrm{out}(R)$, the outputs still
possible, and the state is \emph{determined} when the upper fiber has one element.} \textcolor{newblue}{Determined states stay determined.}

\paragraph{Events.}
\textcolor{newblue}{For a state $(i,R)$ and a legal input $j$, the \emph{input event}}
\begin{equation}
  (i,R)\xrightarrow{\mathsf{input}(j)}
  \bigl(j,\ R[\mathrm{in}\InputGe j]\bigr)
  \label{eq:input-event}
\end{equation}
\textcolor{newblue}{makes $j$ current and \emph{retires} the rows whose input does not
\inrelbase $j$.} \removednote{\textbf{formal definition/notation}, 18 words, rank 63; \textit{We call this input retirement .}} \textcolor{newblue}{Because inputs are cumulative, $\mathsf{input}(j)$ then
$\mathsf{input}(k)$ has the effect of $\mathsf{input}(k)$ alone.}

A system controls what it tells the world. We study what it can safely
promise without assuming that clients tolerate retractions. Later outputs
must therefore \outrelbase earlier ones. \removednote{\textbf{explanation}, 35 words, rank 36; \textit{At state [math], the current input has a unique row [math], so [math] is the only ...}} \removednote{\textbf{formal definition/notation}, 13 words, rank 73; \textit{Let [math] be a state and [math] its row at the current input.}}

\begin{definition}[Output event, seal, commitment]
\label{def:output-event}
\textcolor{newblue}{Let $(i,o)$ be the row at the current input of a state $(i,R)$.} \textcolor{newblue}{The
\emph{output event} $\mathsf{output}(o)$ atomically restricts the remainder
table to the rows whose outputs \outrelbase $o$, and emits $o$:}
\begin{equation}
  (i,R)\xrightarrow{\mathsf{output}(o)}\bigl(i,\ R[\mathrm{out}\OutputGe o]\bigr).
  \label{eq:output-action}
\end{equation}
\textcolor{newblue}{When this removes any row (at some input $j \InputGe i$), we call the state change the \emph{seal} $s_o$: it
\emph{seals} the inputs of the removed rows, and it is a \emph{commitment}.} \removednote{\textbf{formal caution}, 17 words, rank 64; \textit{The current row [math] remains, since [math] itself, so the event always ends at a state.}}
\end{definition}

A commitment removes rows that input alone would not have removed: the inputs
it seals are no longer legal. \removednote{\textbf{example}, 11 words, rank 82; \textit{Seals of this form are the only commitments this section uses.}} \textcolor{newblue}{Systems may commit more than an output
requires, as the registrar does in Figure~\ref{fig:election-frames}(d), but
extra commitments only restrict behavior (Section~\ref{sec:systems}).}

\begin{definition}[Exposable output]
\label{def:exposable}
The output $o$ at $(i,R)$ is \emph{exposable} when $\mathsf{output}(o)$
applies no commitment. \removednote{\textbf{formal definition/notation}, 165 words, rank 8; \textit{In a functional specification this means [math]: the event removes nothing.}} \removednote{\textbf{unclassified}, 27 words, rank 47; \textit{When runs take different paths through the specification, can they still be brought to a common ...}} \removednote{\textbf{formal definition/notation}, 14 words, rank 72; \textit{The transition system consists of specification states and the input and output events above.}}
\end{definition}

\textcolor{newblue}{In a functional specification, $o$ is exposable exactly when $f(j)$ \outrel
$o$ for every legal $j$.}

\paragraph{Runs and consistency.}
\textcolor{newblue}{We now consider the executions corresponding to a specification.}
\begin{definition}[Run]
\label{def:run}
\textcolor{newblue}{A \emph{run} is a finite sequence of input and output events from the
canonical state $\sigma_S(i)$ at some input $i$.} \removednote{\textbf{formal definition/notation}, 32 words, rank 39; \textit{A run is commitment-free when none of its output events applies a commitment.}} \textcolor{newblue}{It is
\emph{commitment-free} when none of its output events applies a commitment, and
a state is \emph{reachable without commitment} when some commitment-free run
reaches it.}
\end{definition}

Input events keep only rows whose input \inrel the current one, and output
events keep only rows whose outputs \outrelbase what has been emitted. \textcolor{newblue}{Both
steps are relational selections, and selections commute; since inputs only \ingrow during
a run, the last input's selection subsumes the earlier ones.}

\begin{lemma}[Rows remaining after a run]
\label{lem:run-residual}
\textcolor{newblue}{If a run ends at current input $i$ after emitting
$o_1,\ldots,o_m$, its remainder table is
$S[\mathrm{in}\InputGe i,\ \mathrm{out}\OutputGe o_1,\ \ldots,\
  \mathrm{out}\OutputGe o_m]$.} \removednote{\textbf{example}, 39 words, rank 33; \textit{the rows whose input [math] and whose output each [math].}; labels eq:run-residual}
\end{lemma}
\textcolor{newblue}{\proofinfull}

To compare runs that end at different inputs, we consider a common \indesc
of those inputs.

\begin{definition}[Consistency]
\label{def:consistency}
An ordered specification is \emph{consistent} when every finite family of
runs is confluent: if the runs end at states
$(i_1,R_1),\ldots,(i_k,R_k)$ and $j$ is a common \indesc of their current
inputs, each run can be continued by further events to one
common state with current input $j$. The runs are then \emph{confluent at} $j$. \removednote{\textbf{explanation}, 16 words, rank 65; \textit{Confluence has a simple test: some row at the common must survive in every run.}}
\end{definition}

\begin{proposition}[Survivor test]
\label{prop:confluence}
\textcolor{newblue}{Runs ending at states $(i_r,R_r)$, $r=1,\ldots,k$, are confluent
at $j$ if and only if some row $(j,u)$ at $j$ is in every $R_r$.} \removednote{\textbf{formal definition/notation}, 0 words, rank 94; labels eq:reconciliation} \textcolor{newblue}{For any such
$u$, $\mathsf{input}(j)$ followed by $\mathsf{output}(u)$ takes every final
state to $\bigl(j,\ S[\mathrm{in}\InputGe j,\ \mathrm{out}\OutputGe u]\bigr)$.}
\end{proposition}

\begin{proof}[Proof sketch]
Rows are never added, so the row of a common state at $j$ was there
in every run. \removednote{\textbf{explanation}, 25 words, rank 51; \textit{Full proof in Appendix .}} \textcolor{newblue}{\fullproof}
\end{proof}

Consistency guarantees runs never diverge irreparably: no ``split brain''. Whatever input
and output events a finite collection of runs has taken, any input \inrelpart all
of theirs leaves a single state that every run can reach: one that keeps every
promise any of them has made and accepts all of that input. \removednote{\textbf{unclassified}, 16 words, rank 66; \textit{Proposition identifies that state.}} \textcolor{newblue}{Consistency is a property of the specification.} \textcolor{newblue}{It describes no protocol and
promises no delivery or termination.}

\subsection{The Functional CALM theorem}
\label{sec:functional-calm-theorem}

Consistency is defined by a specification's executions, monotonicity by its
table and orders alone. Functional CALM shows that these coincide. \textcolor{newblue}{Full proofs of this and later results are in \extended.}

\begin{theorem}[Functional CALM]
\label{thm:functional-calm}
A functional ordered specification is consistent if and only if it is
monotone.
\end{theorem}

\begin{proof}[Proof sketch]
A monotone $f$ only emits outputs that all of its future outputs \outrelbase,
so the row $(j,f(j))$ at a common \indesc $j$ survives in every run. If $f$ is not monotone, a run that emits $f(i)$ and a run that accepts an
\indesc $k$ of $i$ with $f(i)\not\OutputOrd f(k)$ can never be reconciled at $k$. \removednote{\textbf{explanation}, 4 words, rank 85; \textit{Full proof in Appendix .}} \textcolor{newblue}{\fullproof}
\end{proof}

In the equality-order election, $f(xy\cdot)=\bot$ and $f(xyx)=x$ are
incomparable: a run that emits $\bot$ at $xy\cdot$ and one that accepts
$xyx$ are the two runs of Section~\ref{sec:election-story} that can never be reconciled. \removednote{\textbf{explanation}, 16 words, rank 67; \textit{In the functional case, consistency has a strong consequence: outputs leave no trace on the state.}}

\begin{corollary}[State is fixed by input]
\label{cor:state-by-input}
In a consistent functional specification, every run ending at current
input $i$ is in the canonical state $\sigma_S(i)$, whatever outputs it emitted
along the way. \removednote{\textbf{explanation}, 45 words, rank 28; \textit{By monotonicity, the outputs a run emits remove no row at an of [math].}}
\end{corollary}
\textcolor{newblue}{\proofinfull}

\textcolor{newblue}{Knowing the input is therefore enough to know the state, and with it the one
output.} Monotonicity also governs exposability:

\begin{proposition}[Exposure boundary]
\label{prop:functional-exposure-boundary}
A functional ordered specification is monotone if and only if every output at
every state reachable without commitment is exposable.
\end{proposition}

\begin{proof}[Proof sketch]
By Corollary~\ref{cor:state-by-input}, monotone $f$ reaches only canonical
states. At $\sigma_S(i)$, which input alone reaches, $f(i)$ is exposable exactly
when $f$ is monotone at $i$. \removednote{\textbf{explanation}, 21 words, rank 60; \textit{Full proof in Appendix .}} \textcolor{newblue}{\fullproof}
\end{proof}

\textcolor{newblue}{A monotone specification can thus emit every current
output without commitment; if monotonicity fails at $i\InputOrd k$ with
$f(i)\not\OutputOrd f(k)$, then at $\sigma_S(i)$ the event $\mathsf{output}(f(i))$
applies $s_{f(i)}$, sealing $k$.}

\textcolor{newblue}{This is the familiar CALM dichotomy stated at the level of the specification
rather than inside a particular programming language.} The C in CALM has been read
both as consistency and as confluence; Definition~\ref{def:consistency} equates them. Classical CALM fits the pattern with inputs $J$ as sets of facts ordered by
inclusion and output tables $Q(J)$ ordered by inclusion. \textcolor{newblue}{Functional
monotonicity becomes $J\subseteq J'\Rightarrow Q(J)\subseteq Q(J')$,
the standard definition of a monotone query~\cite{abiteboul1995foundations},
which positive Datalog satisfies and negation may not.} \removednote{\textbf{formal definition/notation}, 0 words, rank 95; labels eq:query-monotonicity} Functional CALM
generalizes the statement to arbitrary input and
output orders, whatever formalism expresses the specification. The usual CRDT state specification has the same functional shape,
with received payloads ordered by inclusion and states ordered by the semilattice
order~\cite{shapiro2011crdt}; Sections~\ref{sec:query-instance}
and~\ref{sec:crdt-instance} treat both.


\section{Complete CALM}
\label{sec:complete-calm}

We now extend CALM theory to allow a fiber to contain several outputs. The monotonicity property we want to test as input
\ingrows from $i$ to $j\InputGe i$ must be lifted to sets: can \emph{some}
output in the fiber at $j$ \outrelbase \emph{whatever} outputs runs emit at $i$? In addition to seal commitments on inputs, we now have a second kind of commitment: \emph{choosing}
among permitted outputs. \textcolor{newblue}{Section~\ref{sec:ambiguity-story} introduces the issue by
example.} \textcolor{newblue}{Section~\ref{sec:lifted-monotonicity} defines relational
specifications and Lifted
Monotonicity.} \textcolor{newblue}{Section~\ref{sec:lifted-output-events} lifts
output events and classifies what they remove, and
Section~\ref{sec:complete-calm-theorem} shows that the run-level results of
Section~\ref{sec:executions} carry over and proves Complete CALM.}

\subsection{Output choice by example}
\label{sec:ambiguity-story}

Figure~\ref{fig:ambiguity-frames} contrasts two specifications that permit
several outputs in one fiber. Transaction ordering requires a choice
between incompatible serialization orders. Traceroute observations, by
contrast, can all be \outrelpast by a larger route set.

\begin{figure}[t]
  \centering
  \small
  \setlength{\tabcolsep}{4pt}
  \begin{minipage}[t]{0.36\columnwidth}
    \raggedright
    (f) After $[A,B]$\par\smallskip
    \begin{tabular*}{\linewidth}[t]{@{}l@{\extracolsep{\fill}}l@{}}
      \toprule
      $\mathrm{in}$ & $\mathrm{out}$ \\ \midrule
      $\{A,B\}$ & \sh{$[A,B]$} \\
      \retired{$\{A,B\}$} & \retired{$[B,A]$} \\ \strikerow
      \bottomrule
    \end{tabular*}\par\smallskip
    {\footnotesize no remaining row \outrel the removed one: a commitment}
  \end{minipage}\hfill
  \begin{minipage}[t]{0.6\columnwidth}
    \raggedright
    (g) After $\{s{\to}\texttt{s1}\}$\par\smallskip
    \begin{tabular*}{\linewidth}[t]{@{}l@{\extracolsep{\fill}}l@{}}
      \toprule
      $\mathrm{in}$ & $\mathrm{out}$ \\ \midrule
      $i$ & \sh{$\{s{\to}\texttt{s1}\}$} \\
      \retired{$i$} & \retired{$\{s{\to}a\}$} \\
      \retired{$i$} & \retired{$\{s{\to}a,\ s{\to}a{\to}\texttt{s2}\}$} \\
      $i$ & \sh{$\{s{\to}\texttt{s1},\ s{\to}a\}$} \\
      $i$ & \sh{$\{s{\to}\texttt{s1},\ s{\to}a,\ s{\to}a{\to}\texttt{s2}\}$} \\
      \bottomrule
    \end{tabular*}\par\smallskip
    {\footnotesize each removed row is \outrelpast by a remaining one: an output
    retirement}
  \end{minipage}
  \caption{Two specifications with several rows at one input, each after one output
  event.  (f) Serial orders of two concurrent transactions, under the equality
  order.  (g) Route sets from a source $s$ to a service on \texttt{s1} and
  \texttt{s2}, ordered by inclusion; $s$ reaches \texttt{s1} directly and
  \texttt{s2} through hop $a$.  Legend as in Figure~\ref{fig:election-frames};
  gray alone may also mark an output retirement.}
  \label{fig:ambiguity-frames}
\end{figure}

\paragraph{Output choice and directedness.}
The examples differ in what happens
to alternative outputs when one is emitted. In panel~(f), neither
serialization order \outrel the other, so emitting one requires a
\emph{choice commitment} that removes the other while the input stays legal:
emitting $[A,B]$ chooses it over $[B,A]$. In panel~(g), larger route sets \outrelbase both the
emitted observation and the outputs of the rows it removes. We call such
removal an \emph{output retirement}; every removed output is
\outrelpast by a remaining one, so emitting the route observation loses no information
according to the output order and needs no
commitment: it is \emph{exposable}.

An output event can thus remove rows in three ways: by sealing inputs, as in
Section~\ref{sec:functional-calm}; by retiring outputs, as in panel~(g); or by
choosing among outputs, as in panel~(f). The atomic event
$\mathsf{output}(o)$ does whichever of these $o$ requires, then emits $o$;
sealing and choosing are commitments, retiring is not.

In Figure~\ref{fig:ambiguity-frames}, serialization requires commitment; traceroute avoids it. The difference between them is the order-theoretic concept of a \emph{directed} fiber:
every finite set of outputs in a given fiber has a common \outnoun in that fiber. Directedness at a single input is not sufficient, however;
we must also \outrelbase prior outputs as input \ingrows. \emph{Lifted Monotonicity}
combines the requirements: as input \ingrows from $i$ to $j$, some output
in the fiber at $j$ \outrel whatever outputs runs emit at $i$. Equivalently,
outputs emitted at inputs with a common \indesc $j$ must have a common
\outnoun in the fiber at $j$: the same test as in Section~\ref{sec:election-story}, but
with many outputs per input. Two transaction scheduling runs that emit $[A,B]$ and $[B,A]$ fail it;
two traceroute runs that emit different route sets pass it, since their union
survives in both.

\subsection{Lifted Monotonicity}
\label{sec:relational-specifications}
\label{sec:lifted-monotonicity}

We begin with our most general definition of a specification, which admits
ambiguity of output.
\begin{definition}[Ordered relational specification]
\label{def:relational-specification}
\textcolor{newblue}{An \emph{ordered relational specification} is
$\Spec=(I,\InputOrd,O,\OutputOrd,S)$ with $S\subseteq I\times O$,
where $(I,\InputOrd)$ and $(O,\OutputOrd)$ are partially ordered sets and $S$
is total on $I$.} \removednote{\textbf{formal definition/notation}, 22 words, rank 57; \textit{where [math] and [math] are partially ordered sets and [math] is total on [math]: every input ...}; labels eq:relational-specification} A row $(i,o)\in S$
says that output $o$ is permitted at input $i$. \textcolor{newblue}{The fiber
$\fiber{i}=\{o\in O:(i,o)\in S\}$ may now contain several
outputs; the specification is \emph{ambiguous} when some fiber does.} \removednote{\textbf{formal definition/notation}, 15 words, rank 71; \textit{of Section may now contain several outputs; the specification is ambiguous when some fiber does.}; labels eq:fiber} A
functional specification is the special case $\fiber{i}=\{f(i)\}$. We call
either kind simply an \emph{ordered specification}. \removednote{\textbf{unclassified}, 12 words, rank 76; \textit{Specification states and input events (Sections and ) apply unchanged to relational tables.}} \removednote{\textbf{unclassified}, 12 words, rank 77; \textit{As input from [math] to [math], two things can go wrong.}}
\end{definition}

\textcolor{newblue}{Specification states and input events
(Section~\ref{sec:framework-brief})
apply unchanged to relational tables; Section~\ref{sec:lifted-output-events} lifts
output events.}

\textcolor{newblue}{As input
\ingrows from $i$ to $j$, can some output in the fiber at $j$ \outrelbase
whatever outputs runs emit at $i$?} \textcolor{newblue}{Two things can go wrong.} A
\emph{stranded output} is an output in the fiber at $i$ with no \outnoun in
the fiber at some \indesc $j$. An \emph{incompatible pair} is two outputs in
the fiber at $i$ with no common \outnoun there. We rule out
each, combine them, and restate the result in the form runs need.

First, we rule out stranded outputs.

\begin{definition}[Relational Monotonicity]
\label{def:relational-monotonicity}
An ordered relational specification satisfies \emph{Relational Monotonicity}
when every output keeps some \outnoun as input \ingrows:
        \begin{equation}
  i\InputOrd j,\ o\in\fiber{i}
  \quad\Longrightarrow\quad
  \exists u\in\fiber{j}:\ o\OutputOrd u.
  \label{eq:relational-monotonicity}
        \end{equation}
\end{definition}

For singleton fibers this is functional monotonicity,
Equation~\ref{eq:functional-monotonicity}. It treats each output separately,
so it does not rule out incompatible pairs, which concern several outputs in
one fiber. For a set $F\subseteq O$ of outputs, outputs $o_1,\ldots,o_k$ are
\emph{compatible in $F$} if some $u\in F$ \outrel all of them, and
\emph{incompatible in $F$} otherwise.

Second, we rule out incompatible pairs.

\begin{definition}[Directed fiber, Directed Outputs]
\label{def:fiber-directedness}
A fiber $F\subseteq O$ is \emph{directed} when every nonempty finite
$X\subseteq F$ is compatible in $F$. An ordered relational specification has
\emph{Directed Outputs} when every fiber $\fiber{i}$ is directed.
\end{definition}

A singleton fiber is trivially directed. \textcolor{newblue}{As with common \indescs of inputs, a
common \outnoun need not be
least: directedness requires some upper bound within the fiber for each finite
set, not a least one or a fixed merge operator, so it is weaker than a
join-semilattice.} \removednote{\textbf{example}, 53 words, rank 23; \textit{For example, the fiber [math] under inclusion is directed, since [math] every member, yet [math] and ...}} For example, the fiber
\[
  \{\{a\},\{b\},\{a,b,c\},\{a,b,d\},\{a,b,c,d\}\}
\]
under inclusion is directed, since $\{a,b,c,d\}$ \outrel every member, yet $\{a\}$ and $\{b\}$ have two
incomparable minimal common \outnouns. Thus
the traceroute fiber is directed: the union of any finite set of route sets in
it is again in the fiber and refines each of them. The two transaction orders in
Figure~\ref{fig:ambiguity-frames}(f), by contrast, are incompatible.

\paragraph{Lifted Monotonicity.}
To state both requirements at once, lift $\OutputOrd$ to sets. For nonempty
$A,B\subseteq O$, define
        \begin{equation}
  A\LiftedOrd B
  \;\Longleftrightarrow\;
  \text{every nonempty finite }X\subseteq A\text{ is compatible in }B,
  \label{eq:joint-output-order}
        \end{equation}
the \emph{joint lift} of $\OutputOrd$. We apply it with $A=\fiber{i}$ and
$B=\fiber{j}$ for $i\InputOrd j$. When $X$ is a single output $o$, it asks
that $o$ have a \outnoun in $\fiber{j}$, which rules out stranded outputs. When $X$
holds several outputs, it asks that they have a common \outnoun in $\fiber{j}$; with
$j=i$, this rules out incompatible pairs.

\begin{definition}[Lifted Monotonicity, monotone specification]
\label{def:lifted-monotonicity}
An ordered relational specification satisfies \emph{Lifted Monotonicity}, and
is called \emph{monotone}, when the mapping $i\mapsto\fiber{i}$ is monotone
under the joint lift:
        \begin{equation}
  i\InputOrd j
  \quad\Longrightarrow\quad
  \fiber{i}\LiftedOrd\fiber{j}.
  \label{eq:lifted-monotonicity}
        \end{equation}
\end{definition}

With singleton fibers it is functional monotonicity
(Definition~\ref{def:functional-monotonicity}). \removednote{\textbf{unclassified}, 121 words, rank 12; \textit{As the reading above shows, it implies Relational Monotonicity and Directed Outputs; conversely, for [math] and ...}} \textcolor{newblue}{Lifted
Monotonicity holds exactly when Relational Monotonicity and Directed Outputs
both do.}

\paragraph{The common-\indesc condition.}
Lifted Monotonicity compares an input with its
own expansions. Runs, however, may emit outputs at incomparable inputs, say
$o_1$ at $i_1$ and $o_2$ at $i_2$, and can be reconciled only at a common
\indesc $j$ of both. The following condition asks for exactly that: a common
\outnoun of $o_1$ and $o_2$ in the fiber at $j$.

\begin{definition}[Common-\indesc condition]
\label{def:common-expansion-condition}
An ordered specification satisfies the \emph{common-\indesc condition} when,
for every finite set of rows $(i_1,o_1),\ldots,(i_k,o_k)$ of $S$ and every
common \indesc $j$ of $i_1,\ldots,i_k$, some $u\in S(j)$ \outrel every
$o_r$.
\end{definition}

For the empty set of rows the condition asks only that $S(j)$ be nonempty, which
follows from $S$ being total on $I$.

\begin{proposition}[Common-\indesc form of Lifted Monotonicity]
\label{prop:common-expansion-form}
An ordered relational specification satisfies Lifted Monotonicity if and only
if it satisfies the common-\indesc condition.
\end{proposition}

\begin{proof}[Proof sketch]
Relational Monotonicity carries each $o_r$ to some
\outnoun in the fiber at $j$, and Directed Outputs gives those a common
\outnoun there. \removednote{\textbf{signposting}, 4 words, rank 86; \textit{Full proof in Appendix .}} \textcolor{newblue}{\fullproof}
\end{proof}

\subsection{Lifted output events}
\label{sec:lifted-output-events}

\textcolor{newblue}{Lifted Monotonicity is a property of the table.} \textcolor{newblue}{We now return to execution.} For a state $(i,R)$ and $o\in R(i)$, the output event $\mathsf{output}(o)$
again keeps exactly the rows whose outputs \outrelbase $o$, ending at
$(i,R[\mathrm{out}\OutputGe o])$ as in Equation~\ref{eq:output-action}, and
emits $o$. With several rows at one input, each removed row falls into one of
three classes. \textcolor{newblue}{Consider a removed row $(i', o')$.} If no row at $i'$ \outrel $o$, the row is in $\Sealed_o$,
and the event \emph{seals} $i'$. If a kept row at $i'$ \outrel $o'$,
the row is in $\Retired_o$, and removing it is an \emph{output retirement}. Otherwise the row is in $\Rejected_o$: $i'$ stays legal, but the event
\emph{chooses} against $o'$. The seal $s_o$ removes $\Sealed_o$, and the choice
$c_o$ removes $\Rejected_o$. Both are commitments; output retirement is not. In Figure~\ref{fig:ambiguity-frames}(f), emitting $[A,B]$ puts $[B,A]$ in
$\Rejected_o$. In panel~(g), the two removed route sets are in $\Retired_o$. In the functional case, $\Retired_o$ and $\Rejected_o$ are empty, so this is
the output event of Definition~\ref{def:output-event}. \removednote{\textbf{signposting}, 5 words, rank 84; \textit{Appendix gives the formal definitions.}} \textcolor{newblue}{The formal definitions
of lifted output events are in \extended.}

As before, $o$ is exposable when $\mathsf{output}(o)$ applies no commitment,
that is, when $\Sealed_o$ and $\Rejected_o$ are empty. \textcolor{newblue}{Equivalently, at every legal \indesc $j$ of $i$, $o$ and each
remaining output at $j$ have a common \outnoun at $j$.}

\subsection{The Complete CALM theorem}
\label{sec:complete-calm-theorem}

A lifted output event ends at $(i,R[\mathrm{out}\OutputGe o])$, as in
Equation~\ref{eq:output-action}. Specification states, input events, runs,
commitment-free runs, and consistency are defined as before, and the results
of Section~\ref{sec:executions} used output events only through that
equation. They therefore lift unchanged: the rows remaining after a run
(Lemma~\ref{lem:run-residual}) and the survivor test
(Proposition~\ref{prop:confluence}). With them, consistency
reduces to Lifted Monotonicity.

\begin{theorem}[Complete CALM]
\label{thm:complete-calm}
An ordered specification is consistent if and only if it satisfies Lifted
Monotonicity, equivalently the common-\indesc condition.
\end{theorem}

\begin{proof}[Proof sketch]
If the common-\indesc condition holds, a common \outnoun $u\in S(j)$
of everything emitted survives in every run. If outputs emitted at inputs with a common \indesc $j$ are incompatible in
$S(j)$, the runs that emit them share no row at $j$ and can never be
reconciled there. \removednote{\textbf{signposting}, 4 words, rank 87; \textit{Full proof in Appendix .}} \textcolor{newblue}{\fullproof}
\end{proof}

The name CALM, for Consistency As Logical Monotonicity, tied the idea to a
logic language. The new, more general expansion is Consistency As Lifted Monotonicity. The condition that characterizes consistency is now a
property of a specification: Lifted Monotonicity, which for a function is
ordinary monotonicity. \textcolor{newblue}{Section~\ref{sec:distributed-consistency} shows that
both theorems hold for distributed histories.}

The same condition governs exposability.

\begin{proposition}[Exposure boundary]
\label{prop:exposure-boundary}
An ordered specification satisfies Lifted Monotonicity if and only
if every output at every state reachable without commitment is exposable.
\end{proposition}

\begin{proof}[Proof sketch]
Under the common-\indesc condition, consistency supplies the common \outnoun
the exposure test asks for. Conversely, the exposure test at canonical
states, one row at a time, builds the common \outnoun the condition
asks for. \removednote{\textbf{signposting}, 4 words, rank 88; \textit{Full proof in Appendix .}} \textcolor{newblue}{\fullproof}
\end{proof}

With one output per input, Lifted Monotonicity and the common-\indesc
condition are both monotonicity, so Complete CALM and
Proposition~\ref{prop:exposure-boundary} reduce to Functional CALM and
Proposition~\ref{prop:functional-exposure-boundary}.


\section{Distributed Execution}
\label{sec:distributed-execution}

Consistency compares independent runs. A distributed execution places one run
at each participant, and the only change from our framework's
perspective is that participants can learn about inputs at other participants
via messages. \removednote{\textbf{signposting}, 20 words, rank 61; \textit{Functional CALM and Complete CALM therefore have a natural distributed reading, including a convergence guarantee derived ...}} \textcolor{newblue}{Functional CALM and Complete CALM therefore have a natural
distributed reading, which Section~\ref{sec:histories} makes precise.} We then
let systems make their own commitments, identify coordination with
commitment, and prove Complete CAP.

\subsection{Distributed histories and consistency}
\label{sec:histories}
\label{sec:distributed-consistency}

A finite set $P$ of \emph{participants} each holds a specification state,
starting at a canonical state $\sigma_S(i_p)$, and changes it only by its own
input and output events. A participant learns input directly from the
environment or in a message from another participant. \removednote{\textbf{formal definition/notation}, 29 words, rank 44; \textit{Participants exchange messages by [math] and [math].}} \textcolor{newblue}{A
\emph{message} $m=(j,d)$ carries a cumulative input $j$, or $\varnothing$
if it carries none, and arbitrary application-specific data $d$; one can envision timestamps, version
vectors, state identities, and so on.} \textcolor{newblue}{A receive does not change the specification state; a
participant at $(i,R)$ that receives $(j,d)$ may incorporate $j$ by its own
$\mathsf{input}(k)$ for a legal common \indesc $k$ of $i$ and $j$.} \removednote{\textbf{formal definition/notation}, 119 words, rank 13; \textit{Fix an ordered specification, participants [math], and an initial input [math] for each.}} \textcolor{newblue}{Data can
influence which events occur next, but never directly changes a specification
state.}

\begin{definition}[Distributed history, consistent cut]
\label{def:distributed-history}
\textcolor{newblue}{A \emph{distributed history} $H$ is a finite set of
\emph{history events}, each an occurrence at one participant of an input
event, an output event, $\mathsf{send}(m)$, or $\mathsf{receive}(m)$, with
a total local order on each participant's events, a matching send at another
participant for each receive, and, for each participant $p$, input and output
events that in local order form a defined transition sequence from
$\sigma_S(i_p)$.} A \emph{consistent cut} $C$ is a
subset of history events closed under
happens-before~\cite{lamport1978time,chandy1985distributed}; the standard term
is unrelated to our notion of consistency. \removednote{\textbf{formal definition/notation}, 20 words, rank 62; \textit{A consistent cut contains a prefix of each participant's history events, so by item 3 it ...}; labels def:participant-run}
\end{definition}

\textcolor{newblue}{At a consistent cut $C$, the input and output events of $C$ at participant
$p$, in local order, form its \emph{participant run}.} \removednote{\textbf{example}, 37 words, rank 35; \textit{The endpoint of [math] is [math].}; labels fig:participant-runs} \removednote{\textbf{formal definition/notation}, 39 words, rank 34; \textit{A continuation of a history at a consistent cut [math] is a history [math] that extends ...}} \textcolor{newblue}{The cut is
\emph{confluent at} $j$, for a common \indesc $j$ of the participants'
current inputs, when some continuation of $C$ brings every participant to the
same state $(j,R)$.} \removednote{\textbf{formal definition/notation}, 22 words, rank 58; \textit{History events at different participants are independent, so this holds exactly when the participant runs at ...}; labels def:distributed-consistency} \textcolor{newblue}{A specification has \emph{distributed consistency} when,
for every finite participant set, every consistent cut of every history is
confluent at every common \indesc of the participants' current inputs.}

\begin{proposition}[Consistency and distributed consistency]
\label{prop:distributed-consistency}
An ordered specification is consistent if and only if it has distributed
consistency.
\end{proposition}

\begin{proof}[Proof sketch]
The participant runs at a cut are a family of runs, and any family
of runs can be placed at separate participants that exchange no messages. \removednote{\textbf{signposting}, 12 words, rank 78; \textit{Full proof in Appendix .}} \textcolor{newblue}{\fullproof}
\end{proof}

\textcolor{newblue}{By Complete CALM, then, a specification has distributed consistency exactly
when it satisfies Lifted Monotonicity.} \textcolor{newblue}{For functional specifications,
Corollary~\ref{cor:state-by-input} also gives convergence: participants with
the same current input are in the same state, regardless of message paths,
arrival order, or the outputs they emitted.} \removednote{\textbf{result}, 50 words, rank 26; \textit{In any history of a consistent functional specification, every participant whose current input at a consistent ...}; labels cor:convergence} \textcolor{newblue}{This is the convergence that
strong eventual consistency~\cite{shapiro2011crdt} asks of replicated state,
derived here from the specification rather than from a merge operator.} With several outputs per input this fails even for consistent specifications:
in Figure~\ref{fig:ambiguity-frames}(g), a run that emitted
$\{s{\to}\texttt{s1}\}$ has retired the route sets that omit it, while a run at
the same input that emitted nothing keeps them. Consistency then guarantees
that such replicas can be reconciled, not that they agree. \removednote{\textbf{unclassified}, 99 words, rank 16; \textit{The data [math] of a message is any application state other than input that a participant ...}} \removednote{\textbf{explanation}, 32 words, rank 40; \textit{The histories above use the output events of Sections --, which apply exactly the commitments an ...}}

\subsection{Systems}
\label{sec:systems}

\begin{figure}[t]
  \centering
  \small
  \setlength{\tabcolsep}{5pt}
  \begin{tabular}[t]{@{}l l@{}}
    \multicolumn{2}{@{}l}{(h) 2PL: $A$ locks ahead of $B$}\\
    \toprule
    $\mathrm{in}$ & $\mathrm{out}$ \\ \midrule
    $\{A,B,C\}$ & \sh{$[A,B,C]$} \\
    $\{A,B,C\}$ & \sh{$[A,C,B]$} \\
    $\{A,B,C\}$ & \sh{$[C,A,B]$} \\
    \retired{$\{A,B,C\}$} & \retired{$[B,A,C]$} \\ \strikerow
    \retired{$\{A,B,C\}$} & \retired{$[B,C,A]$} \\ \strikerow
    \retired{$\{A,B,C\}$} & \retired{$[C,B,A]$} \\ \strikerow
    \bottomrule
  \end{tabular}
  \hfill
  \begin{tabular}[t]{@{}l l@{}}
    \multicolumn{2}{@{}l}{(i) MVTO: $B$ receives the first timestamp}\\
    \toprule
    $\mathrm{in}$ & $\mathrm{out}$ \\ \midrule
    $\{A,B,C\}$ & \sh{$[B,A,C]$} \\
    $\{A,B,C\}$ & \sh{$[B,C,A]$} \\
    \retired{$\{A,B,C\}$} & \retired{$[A,B,C]$} \\ \strikerow
    \retired{$\{A,B,C\}$} & \retired{$[A,C,B]$} \\ \strikerow
    \retired{$\{A,B,C\}$} & \retired{$[C,A,B]$} \\ \strikerow
    \retired{$\{A,B,C\}$} & \retired{$[C,B,A]$} \\ \strikerow
    \bottomrule
  \end{tabular}
  \caption{Serial orders of three transactions after one commitment:
  (h) 2PL orders $A$ before $B$; (i) MVTO places $B$ first.  Only complete rows without aborts or reads
  are shown.}
  \label{fig:cc-commitments}
\end{figure}

\textcolor{newblue}{The output events above apply exactly the commitments an output requires.} Many real systems behave differently. \removednote{\textbf{formal definition/notation}, 76 words, rank 19; \textit{For example, the registrar of Section emits [math] and closes the polls in the same event ...}} \textcolor{newblue}{The registrar of
Section~\ref{sec:election-story} closes the polls when it emits $\bot$, and a
lock manager commits to an ordering of ongoing transactions ($A$ before $B$)
without emitting that ordering until later.} We model such systems by letting
them bring bespoke commitments. \removednote{\textbf{unclassified}, 0 words, rank 96; labels def:commitment} \textcolor{newblue}{In general, a \emph{commitment} is a partial
function on specification states that leaves the current input unchanged and
removes a nonempty set of rows, none of them \outrelpast by a remaining row at
the same input.} \removednote{\textbf{result}, 33 words, rank 38; \textit{The seals and choices of Definition and Section satisfy this definition: a seal removes every row ...}; labels eq:commitment-excludes, prop:seal-or-choose} \textcolor{newblue}{At each input where it removes rows, a commitment either
removes every row there, \emph{sealing} that input, or removes some outputs
and keeps others, \emph{choosing} among them.} \removednote{\textbf{formal definition/notation}, 523 words, rank 2; \textit{At an input [math], either no row remains or some does.}; labels def:strategy, def:system, def:system-output-event, eq:lifted-validity} \removednote{\textbf{formal definition/notation}, 26 words, rank 49; \textit{We call the runs of Sections -- canonical runs , and a run under a basis ...}} \textcolor{newblue}{A \emph{system} consists of a
\emph{commitment basis} $\mathcal B$ and a nondeterministic strategy per
participant that picks its next output, send, or receive event from its local
history.} \textcolor{newblue}{Its output event $\Out{\gamma}{o}$ applies a sequence $\gamma$ of
commitments from $\mathcal B$, retires rows, and emits $o$; $\varepsilon$ is
the empty sequence.} \textcolor{newblue}{The full paper gives these definitions formally and shows
that every seal and choice of Section~\ref{sec:complete-calm} is a commitment
in this sense.}

\textcolor{newblue}{Two-phase locking (2PL) and multiversion timestamp
ordering (MVTO)~\cite{bernstein1981concurrency} have different bases: 2PL's
$\mathsf{order}(A,B)$ fixes $A$ before $B$ and nothing else, while MVTO's
$\mathsf{stamp}(B)$ places $B$ before every other transaction
(Figure~\ref{fig:cc-commitments}).} \textcolor{newblue}{Both are choices, but they leave
different alternatives open.} \textcolor{newblue}{Although each
system can define its own basis, no basis yields behavior that the runs of
Sections~\ref{sec:functional-calm}--\ref{sec:complete-calm}, which we call
\emph{canonical runs}, cannot.} \textcolor{newblue}{Every run under any basis has an observably
equivalent canonical run, one that accepts the same inputs, emits the same
outputs, and keeps at least the original's rows after each event (\extended).} \textcolor{newblue}{When 2PL orders $A$ before $B$ while reporting only $A$'s read, the canonical
run reports the read with no commitment; when 2PL later emits a serial order,
so does the canonical run, and both end with the same table.}

\textcolor{newblue}{The canonical commitments are therefore universal.} Locks, timestamps,
quorums, and consensus differ in how and when they decide. \removednote{\textbf{unclassified}, 68 words, rank 20; \textit{Yet every output event of any system removes at least the rows that the canonical event ...}} \textcolor{newblue}{Yet every output
event of any system removes at least the rows that the canonical event for its
output removes.} \textcolor{newblue}{Whether a
specification needs coordination is thus a question about the seals and
choices its outputs require, whatever mechanism a system uses to make them.} \removednote{\textbf{unclassified}, 101 words, rank 15; \textit{Consider 2PL on input [math], where [math] reads a key [math] that no transaction writes, and ...}}

\subsection{Coordination}
\label{sec:coordination}

A commitment is a binding decision that removes behavior the specification
still permits. \textcolor{newblue}{What mutexes, locks, timestamps, and agreement protocols have
in common, semantically, is that each implements commitments, by sealing and
choosing.} A commitment needs
coordination, however, only when some other behavior could contradict it:
another run accepting an input it sealed, or emitting an output it chose
against. We take \emph{coordination} to mean such a commitment, and
\emph{distributed coordination} to mean one that a run at another participant
could contradict. A
mechanism used without committing is not coordination. Closing the polls
after every voter has voted still runs the mechanism, but no ballot is left to
seal, so it removes nothing the specification still permits and applies no
commitment.

\begin{definition}[Coordination]
\label{def:coordination}
Let an output event $\Out{\gamma}{o}$
of a run apply a commitment, and let $\mathcal R$ be a set of runs. The
commitment is \emph{contradicted} in $\mathcal R$ when some run in
$\mathcal R$ accepts an input it seals, or emits, at the same input, an output
it chooses against. A run or history
\emph{coordinates} relative to $\mathcal R$ when one of its output events
applies a contradicted commitment, and is \emph{coordination-free} relative to
$\mathcal R$ otherwise. For a specification, $\mathcal R$ is its canonical
runs; for a system, it is the runs it can perform, under any inputs the
environment may supply. The coordination is \emph{distributed} when the
contradicting run is at another participant of the same history.
\end{definition}

\textcolor{newblue}{For a specification, every commitment is contradicted: some canonical run
accepts each input that a commitment can seal and emits each permitted
output, and that run may be at another participant.} \removednote{\textbf{explanation}, 49 words, rank 27; \textit{Since any family of canonical runs forms a message-free history (Proposition ), that run may be ...}} \textcolor{newblue}{For a specification,
then, coordination-free means commitment-free.} For a system, the answer
depends on which runs it can perform, and here seals and choices differ. \textcolor{newblue}{The
environment may supply any input to any participant, so a seal is distributed
coordination whenever another participant may receive the sealed input.} \removednote{\textbf{example}, 28 words, rank 46; \textit{Closing the polls before every voter has voted seals the remaining ballots, so it is coordination, ...}} Outputs, by contrast, are chosen only by strategies, so a choice is
contradicted only when some run's strategy may emit the output it rules out. \removednote{\textbf{signposting}, 30 words, rank 42; \textit{The next result explains why distributed coordination needs exclusion: a choice and an output it rules ...}}

\begin{proposition}[Conflicting output events cannot be reconciled]
\label{thm:commitments-require-coordination}
From a state $(i,R)$, let one participant perform an output event after which
no output at $i$ \outrel $x$, and let another perform an output event that
emits $x$. The two resulting states cannot be brought to a common state with
current input $i$.
\end{proposition}

\begin{proof}[Proof sketch]
After the second event every output at $i$ \outrel $x$, and after
the first none does; since events never add rows, no row at $i$ can be common
to both. \removednote{\textbf{signposting}, 52 words, rank 25; \textit{Full proof in Appendix .}} \textcolor{newblue}{\fullproof}
\end{proof}

\textcolor{newblue}{The conflict may be transient: an input $j$ \inrelpart $i$ might reconcile the
two outputs.} \textcolor{newblue}{Consistency still counts it as a failure.} \textcolor{newblue}{The environment may
act on what it observes, through channels the system does not see, as in
Lamport's ``anomalous behavior''~\cite{lamport1978time}, and the reconciling
input need not arrive.} \removednote{\textbf{explanation}, 31 words, rank 41; \textit{Consistency therefore excludes every observable conflict, including a transient one.}} \textcolor{newblue}{The system must therefore exclude one of the two
events: by a mutex or atomic instruction in shared memory, by agreement in a
distributed system, or by a rule fixed in advance, such as assigning the
choice to a single participant, which keeps the coordination at that
participant.} \removednote{\textbf{formal caution}, 135 words, rank 11; \textit{Seals are similar, with input in place of output, but the environment may deliver any input ...}} \textcolor{newblue}{In a shared-nothing database in which each participant is the
lock manager for its own keys~\cite{stonebraker1986sharednothing}, ordering
two transactions on one key is coordination but not distributed coordination;
it still needs messages, and a client cut off from the owner cannot get its
answer.} \textcolor{newblue}{Two-phase commit orders a transaction across keys without
distributed coordination: each owner's prepare is its own commitment, and the
coordinator's decision counts votes from a fixed set of participants, which is
monotone.} Communication and distributed coordination therefore differ. \removednote{\textbf{unclassified}, 43 words, rank 30; \textit{Confluence is a safety property, and reaching it is liveness (Section ).}}

\subsection{The Complete CAP Theorem}
\label{sec:complete-cap}

Brewer conjectured that a distributed system cannot be consistent, available,
and partition tolerant at once~\cite{brewer2000towards}; Gilbert and Lynch
proved it for a linearizable register~\cite{gilbert2002cap}. Their theorem
concerns one guarantee, not which guarantees can stay available and
consistent through a partition. Complete CAP answers that question for any
declared guarantee: a specification of what is correct, together with the
rows whose outputs must stay available. The answer is that rows emitted apart
must have a common \outnoun at every common \indesc of their inputs; when every permitted
output must stay available, this is Lifted Monotonicity. A partition does not
change which commitments a specification needs; it removes the ability to make
them jointly, since separated participants cannot agree to seal
an input or make a choice that another could contradict.

\begin{definition}[Partitioned history, Availability, Partition Consistency]
\label{def:partition-notions}
\leavevmode \removednote{\textbf{unclassified}, 970 words, rank 1; \textit{The induced specification is an ordered specification, and the input of an execution only expands as ...}; labels prop:declared-specification}
\begin{enumerate}
\item
\emph{Partitioned history.}  A \emph{partitioned history} divides the
        participants into two or more components and contains no receive whose
        matching send lies in another component; participants still receive
        environment input and emit outputs.
\item
\emph{System.}  Fix a system; we
        require a strategy for each participant of every finite
        participant set and every choice of initial inputs, each depending only
        on its participant's local history.
\item
\emph{Environment.}  The environment is adversarial: it may supply any
        legal input event to any participant.
\item
\emph{Availability.}  A continuation of a partitioned history is
        \emph{local to} a participant $p$ when the system produces it, it
        adds only events of $p$'s component and no input events, and none of
        its receives takes a message sent from another component. Fix an
        \emph{availability relation} $A\subseteq S$, the rows $(i,o)$ whose
        output $o$ must stay available at input $i$, and write $A(i)$ for its
        fiber $\{o:(i,o)\in A\}$. The system has \emph{Availability for $A$}
        when, in every partitioned history it produces, each participant $p$,
        ending in state $(i,R)$, can emit each output in $R(i)\cap A(i)$ in
        some continuation local to $p$. \emph{Availability} is Availability
        for $S$.
\item
\emph{Partition Consistency.}  The system has \emph{Partition
        Consistency} when the final cut of every such history is
        confluent at every common \indesc of the participants' current
        inputs. The
        continuation that confluence asks for need not stay partitioned and
        need not be produced by the system.
\item
\emph{Consistency on $A$.}  $\Spec$ is \emph{consistent on $A$} when
        every finite set of rows of $A$ with a common \indesc $j$ has a common
        \outnoun in $\fiber{j}$.
\end{enumerate}
\end{definition}

With $A=S$, consistency on $A$ is the common-\indesc condition
(Definition~\ref{def:common-expansion-condition}). \removednote{\textbf{formal caution}, 25 words, rank 52; \textit{Unlike Availability, Availability for a smaller [math] can be met by a system whose output events ...}}

\begin{theorem}[Complete CAP]
\label{thm:complete-cap}
For an ordered specification $\Spec$ and an availability relation
$A\subseteq S$, the following are equivalent:
\begin{enumerate}
\item
some system has Availability for $A$ and Partition Consistency;
\item
$\Spec$ is consistent on $A$.
\end{enumerate}
\end{theorem}

\begin{proof}[Proof sketch]
Consider the system whose strategies emit each output of
$A$ by the canonical output event of Section~\ref{sec:complete-calm}. It has
Availability for $A$, and consistency on $A$ gives the outputs it emits a
common \outnoun at every common \indesc of their inputs. Conversely, given one
participant per row of $A$, each in its own component, Availability lets each
emit its row's output, and Partition Consistency then supplies the common
\outnoun. \removednote{\textbf{signposting}, 4 words, rank 89; \textit{Full proof in Appendix .}} \textcolor{newblue}{\fullproof}
\end{proof}

\begin{corollary}
\label{cor:complete-cap-S}
Some system has Availability and Partition Consistency if and only if $\Spec$
satisfies Lifted Monotonicity. When it does, a commitment-free system has
both: one whose strategies may emit any exposable output $o$ by
$\Out{\varepsilon}{o}$. \removednote{\textbf{explanation}, 24 words, rank 54; \textit{With [math], consistency on [math] is Lifted Monotonicity, and the exposure boundary makes every output of ...}}
\end{corollary}
\textcolor{newblue}{\proofinfull}

The converse never assumes that the system is commitment-free, so commitments
cannot enlarge the class of specifications that admit a system with
Availability for $A$ and Partition Consistency. Nor do the commitments of the
canonical system above need coordination. Under consistency on $A$, the canonical event
for an output of $A$ seals nothing and chooses only against outputs outside
$A$, which the system never emits, so it is
coordination-free (Definition~\ref{def:coordination}). For $A=S$ it needs no
commitment at all. \removednote{\textbf{formal caution}, 26 words, rank 50; \textit{For a smaller [math] such local choices can be necessary: if [math] and [math] have no ...}} When consistency on $A$ fails, some
output of $A$ either waits for another component, sacrificing Availability, or
is emitted independently at the cost of Partition Consistency.

\paragraph{Gilbert and Lynch.}
Their theorem is a case of Complete CAP. Consider a register whose input is the
set of operations a participant knows, with their real-time precedence, and
whose outputs answer every operation of the input as some linearization
does~\cite{herlihy1990linearizability}. Let $p$'s input be a write of
$v\neq v_0$ and $q$'s a read; the only outputs they permit are the
acknowledgment and $v_0$. Where the write precedes the read, every output
returns $v$, so $v_0$ has no \outnoun there and $\Spec$ is not consistent on
$S$. The converse in the proof of Theorem~\ref{thm:complete-cap} needs
Availability only so that $p$ can emit the acknowledgment and $q$ can emit
$v_0$. Each is the only output its input permits, so it is enough that every
request receive some response, which is Gilbert and Lynch's availability. It
needs Partition Consistency only to give a common \outnoun of the emitted
outputs, which is what linearizability requires. \textcolor{newblue}{The full paper takes precedence from real time, extends
Complete CAP to systems that choose their responses when responses are not
forced, and derives
the theorem in Gilbert and Lynch's asynchronous model~\cite{hellerstein2026completecalmfull}.}


\section{Composition and Proper Coordination}
\label{sec:composition}

Coordination poses a puzzle. A non-monotone program needs coordination, yet
coordination is itself non-monotone. A key-value store with mutable keys is
non-monotone~\cite{alvaro2011bloom}. State-machine
replication~\cite{schneider1990smr} wraps it in a log. Clients send requests,
the log puts the requested commands in one order, and the key-value store
executes the log and answers each request. The log's specification is also not
monotone: when two commands arrive concurrently, either may be appended next,
and the two resulting logs have no common \outnoun. Yet replicated stores are
consistent, and they coordinate only in the log. Multi-Paxos agrees on each
log entry~\cite{lamport2001paxos}, and the replicas then execute the agreed log
with no further agreement.

That puzzle is explained by our framework. Specifications are
tables, so they compose by equijoin, and composition preserves Lifted
Monotonicity. Because commitments are explicit, we
can also ask which component makes a composition's commitments. When the
upstream component makes every commitment the composition needs, and the
downstream needs none on the inputs the upstream offers, we call the
coordination \emph{proper}
(Section~\ref{sec:proper-coordination}). \removednote{\textbf{signposting}, 4 words, rank 90; \textit{Proofs are in Appendix .}}
\label{sec:sequential-composition}

Intuitively, sequential composition feeds the output of one specification to
the input of the next. The composite is again a table: it pairs an input of
the first specification with every output that the second permits for some
intermediate value that the first permits. Formally, let
$\Spec_1=(I,\InputOrd,M,\preceq_M,S_1)$ and
$\Spec_2=(M,\preceq_M,O,\OutputOrd,S_2)$ be ordered relational specifications
sharing the ordered domain $(M,\preceq_M)$. Their \emph{sequential
composition} has table
\begin{equation}
  S_2\circ S_1
  =\pi_{I,O}(S_1\Join_M S_2)
  =\{(i,o):\exists m.\ (i,m)\in S_1\land(m,o)\in S_2\},
  \label{eq:spec-composition}
\end{equation}
where $\pi$ and $\Join$ are relational projection and equijoin. \textcolor{newblue}{Since $S_1$
is total on $I$ and $S_2$ is total on $M$, every $i\in I$ has a nonempty fiber
in $S_2\circ S_1$.} When the two stages
run as a pipeline, $\Spec_2$ takes the outputs of $\Spec_1$ as its inputs, and
we say that $\Spec_1$ \emph{guards} $\Spec_2$. \removednote{\textbf{signposting}, 22 words, rank 59; \textit{Appendix gives each stage an input of its own as well, as when a store serves ...}}

\begin{lemma}[Composition]
\label{lem:composition}
If $\Spec_1$ and $\Spec_2$ satisfy Lifted Monotonicity, then the ordered
specification with table $S_2\circ S_1$ satisfies Lifted Monotonicity.
\end{lemma}

\begin{proof}[Proof sketch]
Take rows of $S_2\circ S_1$ whose inputs have a common \indesc $j$. The
common-\indesc condition for $\Spec_1$ gives one intermediate value at $j$
that \outrel each row's intermediate value, and that value is a common \indesc
at which the condition for $\Spec_2$ supplies the common \outnoun. \removednote{\textbf{signposting}, 4 words, rank 91; \textit{Full proof in Appendix .}} \textcolor{newblue}{\fullproof}
\end{proof}

By associativity the lemma extends to finite pipelines. \removednote{\textbf{formal definition/notation}, 16 words, rank 68; \textit{For a function [math], its table [math] satisfies Lifted Monotonicity exactly when [math] is monotone (Section ...}} \removednote{\textbf{explanation}, 29 words, rank 45; \textit{The composite is itself an ordered specification.}}

\subsection{Composite CALM}
\label{sec:composite-calm}

\textcolor{newblue}{By Complete CALM, $S_2\circ S_1$ is consistent exactly when it
satisfies Lifted Monotonicity, and then no output needs a commitment.} That test is on the composite's table. When the stages run as a pipeline, we
also want them to avoid commitments. \removednote{\textbf{example}, 238 words, rank 7; \textit{A composite can be jointly monotone in three ways.}} \textcolor{newblue}{A composite can be jointly monotone, with
no stage making a commitment, in three ways.} \textcolor{newblue}{First, every stage is monotone,
as when a monotone client query observes CRDT state
(Section~\ref{sec:crdt-instance}), and Lemma~\ref{lem:composition} applies.} \textcolor{newblue}{Second, the upstream can filter out the inputs on which the downstream is not
monotone: a counter capped at $9$ never produces the counts at which a test
of whether the count is below $10$ answers no.} \textcolor{newblue}{Third, the
order on $M$ is not part of the composite's contract, and a different order
may make two non-monotone stages monotone, as with two stages that each
negate a boolean.} \textcolor{newblue}{These cases do not cover every monotone composite
(\extended).} \textcolor{newblue}{When none applies, some stage may have to commit even though the
composite is monotone, as we discuss next.} \removednote{\textbf{signposting}, 12 words, rank 79; \textit{The replicated store of the opening extends the filtering case of Proposition (2).}} We say that $S_2$ satisfies Lifted
Monotonicity \emph{on} a set $D\subseteq M$ when the common-\indesc condition
holds for its rows with inputs in $D$, at \indescs in $D$. \removednote{\textbf{result}, 102 words, rank 14; \textit{The composite [math] satisfies Lifted Monotonicity, and so is consistent and needs no commitment, in each ...}; labels prop:joint-monotone}

\subsection{Proper Coordination}
\label{sec:proper-coordination}

\textcolor{newblue}{The replicated store of the opening extends the filtering case.} The
key-value store is not monotone
on its own, but it is monotone on what the log offers: a log fixes the order
of commands, and a longer log only extends the responses. Unlike the
filtering case, the upstream must commit to offer only such inputs.

Let $\Spec_1$, which guards $\Spec_2$, be run by a system with any basis, and
let $\Spec_2$ run under a system of its own. Write $\Spec_2|_{\Spec_1}$ for $\Spec_2$ restricted to the outputs of
$\Spec_1$.

\begin{definition}[Proper Coordination]
\label{def:proper-coordination}
A pipeline in which $\Spec_1$ guards $\Spec_2$ is \emph{proper} when
$\Spec_2|_{\Spec_1}$ needs no commitment: every output at every state it
reaches without commitment is exposable.
\end{definition}

\begin{theorem}[Proper Coordination]
\label{thm:proper-coordination}
Fix any system for $\Spec_1$.
\begin{enumerate}
\item
The pipeline is proper exactly when $\Spec_2$ satisfies Lifted
        Monotonicity on the outputs of $\Spec_1$.
\item
If it is proper, $\Spec_2|_{\Spec_1}$ under the empty basis follows
        every run of $\Spec_1$. The result is a run of $S_2\circ S_1$, and
        whenever runs of $\Spec_1$ can be brought to a common state, the
        composite runs are confluent.
\item
If it is proper and $S_2\circ S_1$ violates Lifted Monotonicity,
        $\Spec_1$ cannot emit some output without a commitment.
\end{enumerate}
\end{theorem}

\begin{proof}[Proof sketch]
(1) is Complete CALM and the exposure boundary
(Theorem~\ref{thm:complete-calm}, Proposition~\ref{prop:exposure-boundary})
applied to $\Spec_2|_{\Spec_1}$. (2) Each output event of $\Spec_1$ keeps
only rows above its output, so each input $\Spec_2$ receives lies above the
previous one, and by (1) every answer is exposable. (3) The exposure boundary
of the composite fails at a canonical state, which is the image of a canonical
state of $\Spec_1$. \removednote{\textbf{signposting}, 4 words, rank 92; \textit{Full proof in Appendix .}} \textcolor{newblue}{\fullproof}
\end{proof}

\textcolor{newblue}{The theorem holds whatever basis $\Spec_1$ uses, and a proper pipeline is
itself a guard for a further stage.} Locks in
the key-value store would be pure overhead: by (1), the store restricted to
logs needs no commitment. \removednote{\textbf{signposting}, 43 words, rank 31; \textit{Sequential composition, products of independent stages, and glue stages that copy, merge, or route outputs build ...}} \textcolor{newblue}{The full paper extends these results to
stages with inputs of their own and to acyclic dataflow graphs, checked node
by node much as types are~\cite{hellerstein2026completecalmfull}.}


\section{Related Work and Mappings}
\label{sec:instances}
\label{sec:related}

Complete CALM provides a unified lens for many results across literatures. For the key related ideas, we provide mappings here to our specification model. Each setting declares an input, an output, and their orders before
applying either CALM theorem. \textcolor{newblue}{The full paper~\cite{hellerstein2026completecalmfull}
gives each mapping in more detail, along with deployments, response
obligations, commitment bases for 2PL and MVTO, registers, and further
applications such as $k$-set agreement, renaming, and snapshots.}

\subsection{Query monotonicity}
\label{sec:query-instance}

As previously discussed (\S\ref{sec:functional-calm}),
the CALM Theorem of Ameloot et al.\ proves an operational result for
relational languages on transducer networks: coordination-free evaluation is possible
iff a program is monotone,
under its data-placement and execution assumptions~\cite{ameloot2013relational,ameloot2015weaker}. Functional CALM does not replace the details of that theorem. It abstracts
the result to a property that can be compared with specifications outside
Datalog. \removednote{\textbf{related-work detail}, 150 words, rank 10; \textit{Later refinements of that theorem change the input order or the notion of coordination.}} \textcolor{newblue}{Later refinements change the input order.} \textcolor{newblue}{When transducers know
how data is placed, more queries become coordination-free, among them the
non-monotone win-move query~\cite{zinn2012winmove}, and Ameloot et al.\
characterized the new classes by domain-distinct and domain-disjoint
monotonicity~\cite{ameloot2015weaker,zinn2012weak}.} \textcolor{newblue}{Each is Functional CALM
under an input order in which $I\cup J$ \inrel $I$ only for the $J$ that the
placement policy leaves a node unable to rule out.} The placement model
justifies that order. Functional CALM takes the order as given.

\textcolor{newblue}{Baccaert and Ketsman's \emph{behaviors} map located input facts to sets of
permitted located outputs~\cite{baccaert2023dcbq,baccaert2026spectrum}; under
inclusion, a behavior is a relational specification.} \removednote{\textbf{related-work detail}, 159 words, rank 9; \textit{Under inclusion of outputs, with inputs only by facts at new nodes, a behavior is a ...}} \removednote{\textbf{related-work detail}, 45 words, rank 29; \textit{A CRDT guarantees that replicas receiving the same updates converge to equivalent states, whatever the order ...}} \textcolor{newblue}{Their characterization,
\emph{node monotonicity}, is the common-\indesc condition for two rows at
disjoint networks, and Complete CAP extends it to overlapping networks
(\extended).}

\subsection{CRDTs: monotone queries over convergent state}
\label{sec:crdt-instance}

\textcolor{newblue}{A CRDT guarantees that replicas receiving the same updates converge to
equivalent states, but the Strong Eventual Consistency property it offers
provides no guarantee on what the environment may observe.} Laddad et al.\
proposed monotonicity as the right test for safe observation from CRDT
state~\cite{laddad2022keep}. \removednote{\textbf{example}, 24 words, rank 55; \textit{Their example is a two-phase set, whose state [math] merges by componentwise union.}} We make the proposed test precise with
Functional CALM. Let $(D,\sqsubseteq,\sqcup)$ be the state join-semilattice. The input is the finite set $U$ of immutable payloads received so far, ordered
by inclusion, and the replica's state is $s(U)=s_0\sqcup\bigsqcup U$. \removednote{\textbf{formal definition/notation}, 12 words, rank 80; \textit{Then [math] implies [math], and [math] is the least state both .}} A
client observation $q:D\to(Y,\preceq_Y)$ gives the client-visible
specification $q\circ s$. By Functional CALM, $q\circ s$ is consistent
exactly when $U\subseteq V$ implies $q(s(U))\preceq_Y q(s(V))$. Monotonicity
of $q$ on all of $D$ suffices, but $q$ need only be monotone between reachable
states. \removednote{\textbf{related-work detail}, 53 words, rank 24; \textit{The two-phase query fails the test: a deleted item grows [math] and shrinks the answer.}} \textcolor{newblue}{The query $\mathit{Add}\setminus\mathit{Del}$ over a two-phase set
fails the test: a deleted item grows $U$ and shrinks the answer.} \textcolor{newblue}{CRDT states
form a join-semilattice, but observations need less: Complete CALM requires
only \emph{some} common \outnoun of query outputs, not a least one.}

\subsection{Highly available transactions}
\label{sec:hat-instance}

Consider a transactional specification. An input is the set of begin, read,
write, and commit requests received so far, ordered by inclusion. \removednote{\textbf{related-work detail}, 307 words, rank 4; \textit{An output is a set of public facts, such as [math], [math], and session-order facts, also ...}} \textcolor{newblue}{An output is
a set of public facts, such as $\mathsf{read}(T,x,v)$ and $\mathsf{commit}(T)$,
also ordered by inclusion, and the isolation level or session
guarantee~\cite{bailis2013hat,terry1994session} determines which fact sets are
permitted.} \textcolor{newblue}{Under read committed, two separated replicas that each expose a
read of a committed version and a commit produce facts whose union contains no
dirty read and is itself permitted, a common \outnoun.} \textcolor{newblue}{Under
serializability, if $T_1$ reads $x=0$ and writes $y=1$ on one replica while
$T_2$ reads $y=0$ and writes $x=1$ on the other, no serial order explains
both sets of facts~\cite{adya2000generalized}.} \textcolor{newblue}{Under snapshot isolation, the
first-committer-wins rule lets only one of two concurrent writers of $x$
commit~\cite{berenson1995critique}.} \textcolor{newblue}{When the permitted fact sets are closed
under union, whatever separated replicas expose has a common \outnoun, so the
guarantee stays available (Corollary~\ref{cor:complete-cap-S}).} Serializability and first-committer-wins instead exclude some combinations of
facts that separated replicas can each expose. Emitting either side then
requires a choice (Proposition~\ref{thm:commitments-require-coordination}),
and a partition prevents the replicas from making it together. \removednote{\textbf{formal caution}, 24 words, rank 56; \textit{Eventual convergence of storage does not help.}}

Mahajan et al.\ showed that no consistency semantics stronger than real-time
causal consistency can be enforced in an always-available, one-way convergent
system~\cite{mahajan2011consistency}. \removednote{\textbf{related-work detail}, 98 words, rank 17; \textit{Their availability, that every request is eventually answered, is a response obligation in our terms (Appendix ...}} \removednote{\textbf{related-work detail}, 258 words, rank 6; \textit{An invariant is a predicate that every state must satisfy.}} \textcolor{newblue}{With real-time precedence recorded in
the input, both halves of their result follow from Complete CAP (\extended).}

\subsection{Invariant confluence}
\label{sec:iconfluence-instance}

\textcolor{newblue}{In the model of Bailis et al.~\cite{bailis2014coordination}, a database state
is a set of versions, replicas merge states by union, and a set $T$ of
transactions is \emph{invariant confluent} for an invariant $\mathcal I$ when
the merge of any two reachable states satisfies $\mathcal I$.} \textcolor{newblue}{We construct a
specification $S_{\mathcal I}$ whose inputs are locally valid
\emph{transaction histories}, partial orders of transactions in which one
precedes another when the replica running the second saw the first, ordered
by extension.} \textcolor{newblue}{Outputs are database states, ordered by inclusion, and
$S_{\mathcal I}$ permits at input $h$ the state of any downward-closed
sub-history of $h$ whose merges all satisfy $\mathcal I$.}

\begin{proposition}[Invariant confluence and Lifted Monotonicity]
\label{prop:iconfluence-lm}
$T$ is invariant confluent if and only if $S_{\mathcal I}$ satisfies Lifted
Monotonicity. \removednote{\textbf{signposting}, 2 words, rank 93; \textit{See Appendix .}}
\end{proposition}
\textcolor{newblue}{\proofinfull}

Invariant confluence is therefore not a separate boundary: for the
specification that exposes database states, it is Lifted Monotonicity, and by
Complete CALM it is consistency. \removednote{\textbf{example}, 260 words, rank 5; \textit{The uniqueness constraint of Bailis et al.'s payroll example shows the failure directly.}; labels sec:operational-models} \textcolor{newblue}{In Bailis et al.'s payroll example, two
replicas each assign ID~7 to a new employee; the history in which neither
creation saw the other \inrel both one-creation histories, but no state it
permits \outrel both.}

\subsection{Operational models and scope}
\label{sec:scope}

Automata are a traditional way to model operational behavior
abstractly~\cite{lynch1987io,lynch1996distributed}. Our model keeps only the
inputs and outputs of that behavior and drops the operational details. The
mapping is straightforward: an input's outputs are those of the runs that
consume it. \textcolor{newblue}{Two standard approaches then apply, both used widely and with
success.} \textcolor{newblue}{Simulation and model checking explore bounded fragments of these
runs and exhibit
violations~\cite{misra1986des,flanagan2005dpor,newcombe2015aws,zhou2021foundationdb}.} \textcolor{newblue}{Proofs declare a specification and show, by the standard obligations of
refinement~\cite{abadi1991refinement,lynch1995forward}, that it represents the
machine: every execution's output is dominated by a declared one, and every
declared row is produced by some execution.} \textcolor{newblue}{Such proofs are routinely
mechanized~\cite{hawblitzel2015ironfleet}, and Lifted Monotonicity is then
checked on the declared specification; \extended{} gives the details.} Work
that models systems as state machines therefore maps naturally into our
model, where Complete CALM assesses its consistency directly. \removednote{\textbf{unclassified}, 346 words, rank 3; \textit{We now make the mapping precise, and then explain how it is used in practice.}; labels def:induced-specification, prop:induced-specification}

Our conditions are semantic. \removednote{\textbf{formal caution}, 60 words, rank 22; \textit{Confluence guarantees that a common exists, not that a deployed system reaches it.}} \textcolor{newblue}{Confluence guarantees that
a common \outnoun exists, not that a deployed system reaches it; protocols,
liveness, and cost belong to the operational model in which a specification
is realized.}


\section{Discussion}
\label{sec:discussion}

The computing literature is rich with coordination mechanisms for managing
concurrency both within and across machines. The CALM Conjecture addressed a
question that was still latent at that time: when are these mechanisms required,
and why? The original conjecture built on emerging work in declarative
networking~\cite{loo2009declarative} for intuition. Inspired by Codd's tradition of
separating specification from implementation~\cite{codd1970relational}, declarative
networking showed that the desired effects of protocols could be written as
queries, and subsequently analyzed and optimized into a variety of mechanisms. Could a signature property of declarative languages---monotonicity---be the
desired criterion?

The intuition was sharpened by a related friction in the literature on parallel
query processing. The widely-cited Google MapReduce
paper~\cite{dean2004mapreduce} departed from prior work on parallel query
processing~\cite{wilschut1991dataflow,graefe1990volcano}, conflating data
redistribution (shuffle) with barrier coordination required for a non-monotone
reduce operator. This coupling of unrelated concerns required monotone operators
like join to block in the MapReduce model unnecessarily. Cost models that count
rounds, from BSP supersteps to MapReduce rounds, carry the same
conflation~\cite{valiant1990bsp,karloff2010mrc}.

The CALM Conjecture hypothesized that the database community's more judicious
use of barriers could be generalized: coordination---of any kind---is required
if and only if the program is non-monotone. The CALM Theorem of Ameloot et al.\
proved it. Relational database research provided the intuition and framework for
CALM, but limited the scope to relational languages, which to date have
still not become mainstream beyond query processing (more on this below).

This paper can be seen as moving one level higher in its declarations, from
queries to views. For our analysis, what is important about a specification is
only the single encapsulated view mapping all of a program's base inputs to its
final output. By eliding the query logic entirely, we not only reduce
complexity, we free ourselves from any particular language, providing a
declarative model compatible with any programming framework.

After raising the abstraction, the focus of what to specify also broadens
naturally. Prior work focused on implementing complex programs
declaratively~\cite{loo2005overlays,loo2009declarative,chu2007declarative,alvaro2011bloom}. This was interesting ``declarative engineering'' with concrete benefits in
analysis and optimization. The model here expands the scope, naturally modeling
abstract families of invariants like CAP and isolation levels that have no
specific application or protocol in mind.

Lifting the abstraction improves the power of analysis as well. Prior program
analyses became enmeshed in the details of implementations, however
declarative. Bloom~\cite{alvaro2011bloom} flagged syntactic non-monotonicity in
programs and conservatively injected coordination code in front of it. Yet
coordination code is itself non-monotone, so adding it to a Bloom program only
adds more non-monotone syntax. Jagannathan and colleagues moved the question
from syntax to declared contracts and
invariants~\cite{sivaramakrishnan2015quelea,kaki2018alone}. Their tools pick,
for each operation or transaction, the weakest consistency or isolation level
that the store offers and that still satisfies the declared properties. We
ask what the specification itself requires, independent of any menu of
levels. Because the canonical commitments are universal
(Section~\ref{sec:systems}), the answer holds for every mechanism. Our
composition results need no program logic, and the criteria are exact on
specifications, not conservative. Proper coordination also certifies when a
stage behind a coordinating one needs no coordination of its own.

Underlying these issues is a reassessment of tradition. What is a ``declarative
specification'', and how little does it require to be useful? The question is
timely: declarative specification is seeing a new dawn due to the rise of
automated code generation. Imperative code is increasingly authored by automated
agents, and the ensuing output is often rife with bugs. This has led to
increasing interest in formal specification languages like
Lean~\cite{demoura2021lean4}. Raising the language abstraction helps prevent
agents from getting lost in the incidental complexity of implementation details;
the higher level language enables and complement agents with proofs of
correctness: neurosymbolic computing \jmh{cite Byron Cook}. This is a renewed
opportunity for logic, dataflow or functional programming languages to capture
the attention of developers and their coding agents. Today, when these languages
are used by developers they typically specify protocols and programs at a
significant level of detail, as in the declarative networking papers. This paper
points to the benefits of specification at an even higher level, offering new
opportunities for analyses as in this paper, as well as the potential for formal
program synthesis and optimization.


\section{Conclusion}
\label{sec:conclusion}

A specification abstracts execution details away, and the results here reduce
to a single question: can a participant act on what it knows without being
contradicted by what it learns later? Mutexes, locks, atomic instructions, and
agreement protocols prevent such contradictions by making commitments. Functional CALM answers the question for functions: a functional specification
stays confluent as input arrives exactly when it is monotone, and then no
output needs a seal. This generalizes the classic result from logical fact
sets to arbitrary input and output orders.

Complete CALM extends the answer to specifications that offer several outputs
for one input. The test is Lifted Monotonicity, which combines Relational
Monotonicity with Directed Outputs. When it fails, some reachable output
requires a commitment, and every commitment does one of two jobs: it seals
input or it chooses among outputs. Because commitments are explicit, the model
can also check where a composition coordinates
(Section~\ref{sec:composition}). For example, a totally ordered log and a
key-value store are both non-monotone at their ordinary specifications, yet
when the log guards the store, the log's commitments are all the composition
needs and the store makes none. Multi-Paxos is one way to distribute those
commitments.

Complete CAP is the partition form of the same equivalence. A system can keep
the outputs a guarantee declares available, and consistent, through partitions
exactly when those outputs have a common \outnoun at every common \indesc of their inputs;
when every permitted output must stay available, this is Lifted Monotonicity.

The generality of our specification model opens a range of new questions
regarding the ambiguity of specifications and their commitments. One direction
extends the boolean criteria of this paper to metrics: a Determination
Complexity theory that assesses \emph{how many} commitments are intrinsically
needed to determine a specification. Another is Determination Provenance: a
given output---or its absence---is determined not only by inputs, but also by
the commitments made in a run, which brings formal models of provenance to
concurrent systems. A third direction is to provide general-purpose static
analysis of CALM and proper coordination in a real programming framework.

\begin{acks}
\textcolor{newblue}{Thanks to Conor Power, whose persistent questioning of the CALM canon prompted this work.} \textcolor{newblue}{Much gratitude to Peter Alvaro, Natacha Crooks, Chris Douglas, Tyler Hou, Martin
Kleppmann, Paris Koutris, Edward Lee, Shulu Li, Vincent Liu, David Chu McElroy,
Sarah Morin, and Dan Suciu for helpful feedback.} \textcolor{newblue}{Thanks also to the
anonymous reviewers for frank feedback on earlier drafts of this paper.}

\textcolor{newblue}{Generative AI tools (Claude and ChatGPT) were used as interactive writing and
reviewing assistants during the preparation of this work, including drafting
and revising text, and checking mathematical arguments with Lean proofs.} \textcolor{newblue}{All
content was reviewed, validated, and approved by the author.}
\end{acks}


\bibliographystyle{ACM-Reference-Format}
\bibliography{references,cidr-extra}

\end{document}
